\documentclass[12pt,a4paper,openany]{report}

%% ============================================================
%% ENCODAGE ET LANGUE
%% ============================================================
\usepackage[utf8]{inputenc}
\usepackage[T1]{fontenc}
\usepackage[french]{babel}
\usepackage{lmodern}

%% ============================================================
%% MISE EN PAGE
%% ============================================================
\usepackage[top=2.5cm, bottom=2.5cm, left=3cm, right=2.5cm,
            headheight=15pt]{geometry}
\usepackage{fancyhdr}
\usepackage{setspace}
\usepackage{parskip}
\setlength{\parindent}{0pt}
\setlength{\parskip}{6pt}

%% ============================================================
%% COULEURS
%% ============================================================
\usepackage{xcolor}
\definecolor{cgpblue}{RGB}{0,71,133}
\definecolor{cgplightblue}{RGB}{219,234,254}
\definecolor{cgpgray}{RGB}{107,114,128}
\definecolor{cgplightgray}{RGB}{243,244,246}
\definecolor{cgpgreen}{RGB}{5,150,105}
\definecolor{cgplightgreen}{RGB}{209,250,229}
\definecolor{cgporange}{RGB}{200,80,10}
\definecolor{cgplightorange}{RGB}{255,237,213}
\definecolor{cgpred}{RGB}{180,30,30}
\definecolor{cgplightred}{RGB}{254,226,226}
\definecolor{cgppurple}{RGB}{109,40,217}
\definecolor{cgplightpurple}{RGB}{237,233,254}
\definecolor{cgpplaceholder}{RGB}{200,210,225}
\definecolor{rolebg}{RGB}{245,247,255}

%% ============================================================
%% GRAPHIQUES
%% ============================================================
\usepackage{graphicx}
\graphicspath{{images/}}
\usepackage{tikz}
\usetikzlibrary{shapes.geometric,arrows.meta,positioning,fit,backgrounds}

%% ============================================================
%% BOITES TCOLORBOX
%% ============================================================
\usepackage[most]{tcolorbox}
\tcbuselibrary{breakable,skins}

%% --- Commande screenshot : placeholder figure ---
%% Usage: \screenshot{nom_fichier}{Légende}{Instructions de capture}
\newcommand{\screenshot}[3]{%
  \begin{figure}[htbp]
    \centering
    \begin{tcolorbox}[
      enhanced, width=0.93\textwidth,
      colback=white, colframe=cgpplaceholder,
      title={\small\textbf{[\,Capture d'\'ecran\,] ---}~\texttt{images/#1.png}},
      coltitle=cgpblue, colbacktitle=cgplightblue,
      fonttitle=\small\sffamily, boxrule=1pt,
      left=8pt, right=8pt, top=6pt, bottom=6pt
    ]
      {\small\textbf{Que capturer :}~#3}\\[8pt]
      \begin{center}
        \begin{tikzpicture}
          \fill[cgplightgray, rounded corners=4pt] (0,0) rectangle (13,4.5);
          \draw[cgpplaceholder, thick, dashed, rounded corners=4pt]
                (0,0) rectangle (13,4.5);
          \node[font=\large\itshape, text=cgpgray] at (6.5,2.25) {%
            \textit{[~#2~]}%
          };
          \node[font=\tiny\sffamily, text=cgpplaceholder] at (6.5,0.4)
                {Remplacer par la capture r\'eelle};
        \end{tikzpicture}
      \end{center}
    \end{tcolorbox}
    \caption{#2}
    \label{fig:#1}
  \end{figure}
}

%% --- Note informative (bleue) ---
\newtcolorbox{notebox}[1][Note]{
  enhanced, breakable,
  title={\small\textbf{\faInfoCircle\enskip #1}},
  colback=cgplightblue, colframe=cgpblue,
  fonttitle=\small\sffamily, coltitle=cgpblue,
  colbacktitle=cgplightblue, boxrule=1pt,
  left=8pt, right=8pt, top=4pt, bottom=4pt
}

%% --- Attention (orange) ---
\newtcolorbox{warnbox}[1][Attention]{
  enhanced, breakable,
  title={\small\textbf{\faExclamationTriangle\enskip #1}},
  colback=cgplightorange, colframe=cgporange,
  fonttitle=\small\sffamily, coltitle=white,
  colbacktitle=cgporange, boxrule=1pt,
  left=8pt, right=8pt, top=4pt, bottom=4pt
}

%% --- Conseil (vert) ---
\newtcolorbox{tipbox}[1][Conseil]{
  enhanced, breakable,
  title={\small\textbf{\faLightbulb\enskip #1}},
  colback=cgplightgreen, colframe=cgpgreen,
  fonttitle=\small\sffamily, coltitle=white,
  colbacktitle=cgpgreen, boxrule=1pt,
  left=8pt, right=8pt, top=4pt, bottom=4pt
}

%% --- Rôle(s) concerné(s) ---
\newtcolorbox{rolebox}[1]{
  enhanced, breakable,
  title={\small\textbf{\faUserShield\enskip R\^ole(s) concern\'e(s) :~#1}},
  colback=rolebg, colframe=cgpblue!50,
  fonttitle=\small\sffamily, coltitle=cgpblue,
  colbacktitle=rolebg, boxrule=0.8pt,
  left=8pt, right=8pt, top=4pt, bottom=4pt
}

%% --- Étapes numérotées ---
\newcounter{stepnum}
\newenvironment{steps}{%
  \setcounter{stepnum}{0}%
  \begin{list}{%
    \refstepcounter{stepnum}%
    \textbf{\color{cgpblue}\'{E}tape~\thestepnum.}%
  }{%
    \setlength{\leftmargin}{2.5cm}%
    \setlength{\labelwidth}{2.2cm}%
    \setlength{\labelsep}{0.3cm}%
    \setlength{\itemsep}{5pt}%
    \setlength{\parsep}{0pt}%
  }%
}{%
  \end{list}%
}

%% ============================================================
%% TABLEAUX
%% ============================================================
\usepackage{booktabs}
\usepackage{longtable}
\usepackage{tabularx}
\usepackage{array}
\usepackage{multirow}
\usepackage{colortbl}
\newcolumntype{L}[1]{>{\raggedright\arraybackslash}p{#1}}
\newcolumntype{C}[1]{>{\centering\arraybackslash}p{#1}}

%% ============================================================
%% LISTES
%% ============================================================
\usepackage{enumitem}
\setlist[itemize]{noitemsep, topsep=3pt, leftmargin=1.5em}
\setlist[enumerate]{noitemsep, topsep=3pt, leftmargin=2em}

%% ============================================================
%% ICÔNES
%% ============================================================
\usepackage{fontawesome5}

%% ============================================================
%% LIENS
%% ============================================================
\usepackage{hyperref}
\hypersetup{
  colorlinks=true, linkcolor=cgpblue,
  urlcolor=cgpblue, citecolor=cgpblue,
  bookmarksopen=true,
  pdftitle={Manuel Utilisateur -- CGP Annuaire Formation},
  pdfauthor={\'Equipe CGP},
  pdfsubject={Documentation utilisateur},
}

%% ============================================================
%% FORMATAGE DES TITRES
%% ============================================================
\usepackage{titlesec}
\titleformat{\chapter}[display]{%
  \normalfont\huge\bfseries\color{cgpblue}%
}{\chaptertitlename~\thechapter}{16pt}{\Huge}
\titlespacing*{\chapter}{0pt}{-20pt}{18pt}

\titleformat{\section}{%
  \normalfont\Large\bfseries\color{cgpblue}%
}{\thesection}{1em}{}
\titlespacing*{\section}{0pt}{14pt}{6pt}

\titleformat{\subsection}{%
  \normalfont\large\bfseries\color{cgpblue!80}%
}{\thesubsection}{1em}{}
\titlespacing*{\subsection}{0pt}{10pt}{4pt}

\titleformat{\subsubsection}{%
  \normalfont\normalsize\bfseries\color{cgpblue!60}%
}{\thesubsubsection}{1em}{}
\titlespacing*{\subsubsection}{0pt}{8pt}{3pt}

%% ============================================================
%% EN-TETES ET PIEDS DE PAGE
%% ============================================================
\pagestyle{fancy}
\fancyhf{}
\fancyhead[L]{\small\textcolor{cgpgray}{\nouppercase{\leftmark}}}
\fancyhead[R]{\small\textcolor{cgpgray}{CGP -- Manuel Utilisateur}}
\fancyfoot[C]{\small\textcolor{cgpgray}{\thepage}}
\fancyfoot[R]{\small\textcolor{cgpgray}{v1.0 -- Avril 2025}}
\renewcommand{\headrulewidth}{0.3pt}
\renewcommand{\footrulewidth}{0.3pt}
\fancypagestyle{plain}{%
  \fancyhf{}%
  \fancyfoot[C]{\small\textcolor{cgpgray}{\thepage}}%
  \renewcommand{\headrulewidth}{0pt}%
  \renewcommand{\footrulewidth}{0.3pt}%
}

%% ============================================================
%% DOCUMENT
%% ============================================================
\begin{document}

%%--------------------------------------------------------------
%% PAGE DE TITRE
%%--------------------------------------------------------------
\begin{titlepage}
\centering
\vspace*{1.5cm}

{\color{cgpblue}\rule{\linewidth}{3pt}}
\vspace{0.6cm}

\begin{tcolorbox}[
  enhanced, width=0.8\textwidth,
  colback=cgplightblue, colframe=cgpblue,
  boxrule=2pt, arc=8pt,
  left=12pt, right=12pt, top=14pt, bottom=14pt
]
  \centering
  {\color{cgpblue}\fontsize{48}{52}\selectfont\bfseries CGP}\\[0.4cm]
  {\color{cgpgray}\Large Annuaire Formation}\\[0.2cm]
  {\color{cgpgray}\normalsize Syst\`eme de Gestion de l'Annuaire Universitaire}
\end{tcolorbox}

\vspace{0.8cm}
{\color{cgpblue}\rule{0.5\linewidth}{1.5pt}}
\vspace{0.7cm}

{\fontsize{32}{36}\selectfont\bfseries\color{cgpblue} Manuel Utilisateur}\\[0.4cm]
{\Large\color{cgpgray} Guide complet d'utilisation --- Tous profils}

\vspace{1cm}

\begin{tcolorbox}[
  enhanced, width=0.55\textwidth,
  colback=cgplightgray, colframe=cgpgray!40,
  boxrule=0.5pt, arc=4pt,
  left=10pt, right=10pt, top=8pt, bottom=8pt
]
\centering\small
\begin{tabular}{@{}ll@{}}
  \textbf{Version :}        & 1.0 \\[4pt]
  \textbf{Date :}           & Avril 2025 \\[4pt]
  \textbf{Audience :}       & Tous les utilisateurs \\[4pt]
  \textbf{Classification :} & Usage interne \\[4pt]
\end{tabular}
\end{tcolorbox}

\vfill
{\color{cgpblue}\rule{\linewidth}{3pt}}
\end{titlepage}

%%--------------------------------------------------------------
%% AVANT-PROPOS
%%--------------------------------------------------------------
\chapter*{Avant-propos}
\addcontentsline{toc}{chapter}{Avant-propos}

Ce manuel d'utilisation d\'ecrit l'ensemble des fonctionnalit\'es de l'application \textbf{CGP -- Annuaire Formation}, un syst\`eme de gestion de l'annuaire universitaire destin\'e \`a l'ensemble des personnels de l'\'etablissement.

\medskip
L'application permet de g\'erer les structures acad\'emiques, les responsables p\'edagogiques, les affectations de r\^oles, les organigrammes et les \'echanges de donn\'ees dans un environnement multi-ann\'ees acad\'emiques.

\medskip
\textbf{\`A qui s'adresse ce manuel ?}

Ce document s'adresse \`a l'ensemble des utilisateurs de la plateforme, quels que soient leurs r\^oles. Un chapitre d\'edi\'e \`a chaque profil principal d\'etaille les fonctionnalit\'es accessibles et les proc\'edures \`a suivre. Les fonctionnalit\'es communes \`a tous les profils sont regroup\'ees dans le chapitre~\ref{chap:commun}.

\begin{warnbox}[Images \`a ins\'erer]
Les encadr\'es gris\'es portant la mention \textbf{[Capture d'\'ecran]} sont des emplacements r\'eserv\'es aux captures d'\'ecran de l'application r\'eelle. Chaque bloc pr\'ecise le nom du fichier image attendu (dans le r\'epertoire \texttt{images/}) ainsi que des instructions d\'etaill\'ees pour r\'ealiser la capture correspondante. Une fois les captures disponibles, remplacer la commande \verb|\screenshot{...}{...}{...}| par une commande \verb|\includegraphics| standard dans un environnement \texttt{figure}.
\end{warnbox}

\begin{tipbox}[Convention de remplacement d'image]
Pour remplacer un placeholder par une vraie image\,:
\begin{verbatim}
\begin{figure}[htbp]
  \centering
  \includegraphics[width=0.9\textwidth]{images/nom_fichier.png}
  \caption{Légende}
  \label{fig:nom_fichier}
\end{figure}
\end{verbatim}
\end{tipbox}

%%--------------------------------------------------------------
%% TABLE DES MATIÈRES
%%--------------------------------------------------------------
\tableofcontents
\clearpage

%%--------------------------------------------------------------
%% CHAPITRES
%%--------------------------------------------------------------
\chapter{Introduction}
\label{chap:introduction}

\section{Pr\'esentation de l'application}

\textbf{CGP -- Annuaire Formation} est une application web de gestion de l'annuaire universitaire. Elle centralise les informations sur les structures acad\'emiques, les responsables p\'edagogiques, les formations et le personnel de secr\'etariat de l'\'etablissement.

Ses principales fonctionnalit\'es sont\,:

\begin{itemize}
  \item \textbf{Annuaire et recherche} : consultation multi-crit\`eres des responsables, formations, structures et secr\'etariats.
  \item \textbf{Gestion des affectations} : association de personnes \`a des r\^oles au sein des structures.
  \item \textbf{Organigrammes} : g\'en\'eration et consultation d'organigrammes hi\'erarchiques (vue structures ou vue personnes).
  \item \textbf{Gestion multi-ann\'ees} : cr\'eation, clonage, archivage et suppression d'ann\'ees acad\'emiques.
  \item \textbf{Import et export} : \'echanges de donn\'ees via des classeurs Excel standardis\'es.
  \item \textbf{D\'el\'egations} : transfert temporaire de droits entre utilisateurs.
  \item \textbf{Signalements} : d\'eclaration et suivi d'erreurs sur les donn\'ees.
  \item \textbf{Journal d'audit} : tra\c{c}abilit\'e compl\`ete des actions effectu\'ees dans l'application.
\end{itemize}

\screenshot{img_accueil}{Page d'accueil de l'application CGP}{Capturer la page d'accueil ou de connexion de l'application telle qu'elle appara\^it dans le navigateur web (Google Chrome ou Firefox recommand\'e). L'URL de l'application doit \^etre visible dans la barre d'adresse. Utiliser une fen\^etre en plein \'ecran, r\'esolution 1920$\times$1080.}

\section{Architecture g\'en\'erale}

L'application est organis\'ee autour de concepts fondamentaux\,:

\begin{description}[style=nextline, leftmargin=3cm, font=\bfseries]
  \item[Ann\'ee acad\'emique] Unit\'e de temps principale. Toutes les donn\'ees (structures, affectations, organigrammes) sont rattach\'ees \`a une ann\'ee. Une seule ann\'ee peut \^etre \og{} EN\_COURS \fg{} \`a la fois.
  \item[Structure (entit\'e)] Unit\'e organisationnelle : Composante, D\'epartement, Mention, Parcours, Niveau, Service, etc.
  \item[Utilisateur] Toute personne ayant un acc\`es \`a l'application.
  \item[Affectation] Lien entre un utilisateur, un r\^ole et une structure pour une ann\'ee donn\'ee.
  \item[R\^ole] Fonction occupée par un utilisateur dans une structure (ex. : Directeur de composante, Responsable de formation).
  \item[Organigramme] Repr\'esentation graphique hi\'erarchique g\'en\'er\'ee \`a partir des affectations.
  \item[D\'el\'egation] Transfert temporaire de droits d'un utilisateur vers un autre.
  \item[Signalement] Rapport d'erreur soumis par n'importe quel utilisateur.
\end{description}

\section{Acc\`es \`a l'application}

L'application CGP est accessible via un navigateur web \`a l'adresse communiqu\'ee par votre administrateur syst\`eme.

\begin{notebox}[Navigateurs compatibles]
L'application est optimis\'ee pour\,: \textbf{Google Chrome} (recommand\'e), \textbf{Mozilla Firefox} et \textbf{Microsoft Edge}. \\
Internet Explorer et les navigateurs anciens ne sont \textbf{pas} support\'es.
\end{notebox}

\begin{notebox}[Pr\'erequis r\'eseau]
  Une connexion au r\'eseau de l'universit\'e est n\'ecessaire (connexion filaire, Wi-Fi universit\'e ou VPN selon la configuration de votre \'etablissement).
\end{notebox}

\section{Organisation du manuel}

Ce manuel est structur\'e comme suit\,:

\medskip
\begin{tabularx}{\textwidth}{C{2.5cm}L{4cm}X}
  \toprule
  \rowcolor{cgplightblue}
  \textbf{Chapitre} & \textbf{Titre} & \textbf{Contenu} \\
  \midrule
  Chapitre~1 & Introduction & Ce chapitre -- pr\'esentation g\'en\'erale \\
  \rowcolor{cgplightgray}
  Chapitre~2 & Connexion et profil & Authentification, profil utilisateur, notifications \\
  Chapitre~3 & R\^oles et droits & Pr\'esentation de tous les r\^oles, matrice des droits \\
  \rowcolor{cgplightgray}
  Chapitre~4 & Services Centraux & Manuel complet pour le r\^ole \texttt{services-centraux} \\
  Chapitre~5 & Directeur de Composante & Manuel pour le r\^ole \texttt{directeur-composante} \\
  \rowcolor{cgplightgray}
  Chapitre~6 & Directions et Responsables & Manuel pour les r\^oles de direction (d\'epartement, mention, etc.) \\
  Chapitre~7 & Secr\'etariat P\'edagogique & Manuel pour \texttt{secretariat-pedagogique} \\
  \rowcolor{cgplightgray}
  Chapitre~8 & Utilisateur Standard & Manuel pour \texttt{utilisateur-simple} et \texttt{lecture-seule} \\
  Chapitre~9 & Fonctionnalit\'es communes & Annuaire, profil, signalements, organigrammes (lecture) \\
  \bottomrule
\end{tabularx}

\medskip
\begin{tipbox}[Comment utiliser ce manuel ?]
  Si vous cherchez la proc\'edure pour une fonctionnalit\'e sp\'ecifique, consultez d'abord le chapitre correspondant \`a votre r\^ole. Si la fonctionnalit\'e n'y est pas d\'etaill\'ee, reportez-vous au chapitre~\ref{chap:commun} (Fonctionnalit\'es communes).
\end{tipbox}

\section{Conventions typographiques}

Les conventions utilis\'ees dans ce manuel sont les suivantes\,:

\medskip
\begin{tabularx}{\textwidth}{C{3cm}X}
  \toprule
  \rowcolor{cgplightblue}
  \textbf{Convention} & \textbf{Signification} \\
  \midrule
  \textbf{Gras} & \'El\'ement d'interface : bouton, menu, onglet (ex. : \textbf{Cr\'eer}, \textbf{Valider}) \\
  \rowcolor{cgplightgray}
  \textit{Italique} & Nom de champ de formulaire (ex. : \textit{Nom}, \textit{Courriel institutionnel}) \\
  \texttt{monospace} & Valeur technique, code, nom de r\^ole (ex. : \texttt{services-centraux}) \\
  \rowcolor{cgplightgray}
  \faInfoCircle & Note informative compl\'ementaire \\
  \faExclamationTriangle & Avertissement -- point important \`a ne pas n\'egliger \\
  \rowcolor{cgplightgray}
  \faLightbulb & Conseil -- bonne pratique recommand\'ee \\
  \faUserShield & Indication du ou des r\^oles concern\'es par la section \\
  \bottomrule
\end{tabularx}

\chapter{Connexion et Profil Utilisateur}
\label{chap:connexion}

\begin{rolebox}{Tous les utilisateurs}
  Ce chapitre s'applique \`a l'ensemble des profils. Chaque utilisateur doit se connecter avant d'acc\'eder \`a l'application.
\end{rolebox}

\section{Page de connexion}

\`A l'ouverture de l'application, l'utilisateur arrive sur la page de connexion. Aucun acc\`es aux donn\'ees n'est possible sans authentification.

\screenshot{img_login}{Page de connexion}{Capturer la page de connexion de l'application dans son int\'egralit\'e. Le formulaire avec le champ identifiant et le bouton de validation doit \^etre clairement visible. Fen\^etre en plein \'ecran, fond de page visible.}

\section{Proc\'edure de connexion}

\begin{steps}
  \item Ouvrir le navigateur web et saisir l'adresse de l'application fournie par l'administrateur.
  \item Sur la page de connexion, saisir votre \textbf{identifiant} (login) dans le champ pr\'evu.
  \item Cliquer sur le bouton \textbf{Se connecter} (ou appuyer sur la touche \textsf{Entr\'ee}).
  \item L'application v\'erifie votre identit\'e et charge votre profil, vos affectations et l'ann\'ee acad\'emique courante.
  \item Vous \^etes redirig\'e automatiquement vers le \textbf{Tableau de bord}.
\end{steps}

\begin{notebox}[Syst\`eme d'authentification]
  L'application utilise un syst\`eme de connexion par \textbf{identifiant unique} (login). Selon la configuration de votre \'etablissement, cet identifiant peut correspondre \`a votre identifiant institutionnel (num\'ero de personnel, alias de messagerie, etc.). Rapprochez-vous de votre administrateur en cas de probl\`eme de connexion.
\end{notebox}

\begin{warnbox}[Identifiant incorrect]
  Si votre identifiant n'est pas reconnu, l'application affiche un message d'erreur. V\'erifiez que\,: (1) vous avez saisi l'identifiant correct, (2) votre compte a bien \'et\'e cr\'e\'e par un administrateur, (3) votre compte poss\`ede au moins une affectation active dans l'ann\'ee EN\_COURS.
\end{warnbox}

\section{Tableau de bord apr\`es connexion}

\`A la connexion, l'application affiche le tableau de bord. Celui-ci pr\'esente\,:

\begin{itemize}
  \item \textbf{L'ann\'ee acad\'emique active} visible dans l'en-t\^ete (ex. : \og 2024-2025 \fg).
  \item Les \textbf{raccourcis} vers les principales fonctionnalit\'es accessibles selon votre r\^ole.
  \item Un \textbf{indicateur de notifications} (cloche) signalant les \'ev\`enements en attente.
  \item Pour les \textbf{Services Centraux uniquement} : un s\'electeur d'ann\'ee permettant de changer de contexte.
\end{itemize}

\screenshot{img_dashboard_general}{Tableau de bord apr\`es connexion}{Capturer la page du tableau de bord telle qu'elle appara\^it apr\`es connexion avec un compte de test. Montrer l'en-t\^ete (ann\'ee courante, cloche de notifications), le menu de navigation lat\'eral et les cartes/tuiles du tableau de bord. Plein \'ecran.}

\section{S\'electeur d'ann\'ee acad\'emique}

\begin{rolebox}{services-centraux uniquement}
  Le changement d'ann\'ee est r\'eserv\'e au r\^ole \texttt{services-centraux}.
\end{rolebox}

Pour les utilisateurs ayant le r\^ole \texttt{services-centraux}, un menu d\'eroulant dans l'en-t\^ete permet de s\'electionner n'importe quelle ann\'ee (EN\_COURS, PREPARATION, ARCHIVEE).

\begin{steps}
  \item Localiser le s\'electeur d'ann\'ee dans l'en-t\^ete de l'application (g\'en\'eralement en haut \`a droite).
  \item Cliquer sur l'ann\'ee actuellement affich\'ee pour d\'eplier la liste des ann\'ees disponibles.
  \item Cliquer sur l'ann\'ee souhait\'ee.
  \item L'application recharge automatiquement toutes les donn\'ees dans le contexte de l'ann\'ee s\'electionn\'ee.
\end{steps}

\screenshot{img_selecteur_annee}{S\'electeur d'ann\'ee (Services Centraux)}{Capturer l'en-t\^ete de l'application avec le menu d\'eroulant du s\'electeur d'ann\'ee ouvert, montrant plusieurs ann\'ees disponibles. Utiliser un compte services-centraux.}

\begin{notebox}
  Pour tous les autres r\^oles, l'ann\'ee affich\'ee est l'ann\'ee \textbf{EN\_COURS} et ne peut pas \^etre modifi\'ee. Elle est visible mais non interactive.
\end{notebox}

\section{Navigation dans l'application}

L'application dispose d'un menu de navigation (lat\'eral ou en haut de page) dont le contenu varie selon le r\^ole de l'utilisateur.

Les rubriques possibles sont\,:

\medskip
\begin{tabularx}{\textwidth}{L{4cm}L{5cm}X}
  \toprule
  \rowcolor{cgplightblue}
  \textbf{Rubrique} & \textbf{Ic\^one / Label} & \textbf{Accessible par} \\
  \midrule
  Tableau de bord   & \faHome\ Tableau de bord  & Tous \\
  \rowcolor{cgplightgray}
  Rechercher        & \faSearch\ Rechercher      & Tous \\
  Organigramme      & \faProjectDiagram\ Organigramme & Tous \\
  \rowcolor{cgplightgray}
  Import / Export   & \faFileExcel\ Import/Export & SC, Directeurs composante/admin \\
  Responsables      & \faUsers\ Responsables    & SC, Directeurs, Admins \\
  \rowcolor{cgplightgray}
  Structures        & \faBuilding\ Structures   & SC, Directeurs composante \\
  D\'emandes de r\^oles & \faClipboardList\ D\'emandes & SC, Directions \\
  \rowcolor{cgplightgray}
  D\'el\'egations   & \faUserCheck\ D\'el\'egations & SC, Directeurs \\
  Ann\'ees          & \faCalendar\ Ann\'ees     & SC uniquement \\
  \rowcolor{cgplightgray}
  Audit             & \faHistory\ Audit         & SC uniquement \\
  Signalements      & \faFlag\ Signalements     & Tous \\
  \rowcolor{cgplightgray}
  Mon profil        & \faUser\ Profil           & Tous \\
  \bottomrule
\end{tabularx}

\medskip
\begin{notebox}
  Les rubriques non accessibles pour un r\^ole donn\'e ne s'affichent pas dans le menu. Le menu est donc adapt\'e automatiquement.
\end{notebox}

\screenshot{img_menu_navigation}{Menu de navigation (exemple Services Centraux)}{Capturer le menu de navigation complet avec toutes les rubriques visibles pour un compte services-centraux (c'est le menu le plus complet). Montrer le menu lat\'eral enti\`erement d\'eploy\'e.}

\section{D\'econnexion}

Pour quitter l'application, cliquer sur l'ic\^one de profil ou le menu utilisateur (g\'en\'eralement en haut \`a droite de l'\'ecran), puis s\'electionner \textbf{D\'econnecter} ou \textbf{Se d\'econnecter}.

\begin{warnbox}[Session navigateur]
  L'application conserve votre identifiant de session dans le navigateur (\texttt{localStorage}). Si vous partagez votre poste de travail, pensez \`a vous d\'econnecter explicitement apr\`es chaque utilisation.
\end{warnbox}

\section{Gestion du profil}

La rubrique \textbf{Mon Profil} (accessible depuis le menu de navigation ou l'ic\^one utilisateur) vous permet de consulter et de modifier une partie de vos informations personnelles.

\subsection{Consultation du profil}

La fiche profil affiche\,:
\begin{itemize}
  \item \textbf{Informations d'identit\'e} : nom, pr\'enom (non modifiables par l'utilisateur).
  \item \textbf{Identifiant de connexion} (non modifiable).
  \item \textbf{R\^oles et affectations} : liste des r\^oles actifs pour l'ann\'ee courante.
  \item \textbf{Coordonn\'ees modifiables} : courriel, t\'el\'ephone, bureau.
\end{itemize}

\screenshot{img_profil}{Page Mon Profil}{Capturer la page profil utilisateur enti\`ere. Montrer les sections informations d'identit\'e, les affectations actives et les champs modifiables (courriel, t\'el\'ephone, bureau). Utiliser des donn\'ees de test non confidentielles.}

\subsection{Modification du profil}

\begin{steps}
  \item Naviguer vers \textbf{Mon Profil} dans le menu.
  \item Dans la section \og Mes coordonn\'ees \fg, cliquer sur \textbf{Modifier}.
  \item Mettre \`a jour les champs souhait\'es\,: \textit{Courriel secondaire}, \textit{T\'el\'ephone}, \textit{Bureau}.
  \item Cliquer sur \textbf{Enregistrer}.
\end{steps}

\begin{warnbox}[Champs non modifiables]
  Certains champs ne peuvent \^etre modifi\'es que par un administrateur ou les Services Centraux\,: \textit{Nom}, \textit{Pr\'enom}, \textit{Identifiant}, \textit{R\^oles/Affectations}. Si une de ces informations est erron\'ee, cr\'eez un signalement (voir chapitre~\ref{chap:commun}).
\end{warnbox}

\section{Notifications}

La cloche \faBell\ en haut de page indique le nombre de notifications non lues. Les notifications informent l'utilisateur des \'ev\`enements le concernant\,:

\begin{itemize}
  \item R\'eponse \`a une demande de r\^ole (approuv\'ee ou refus\'ee).
  \item Signalement cl\^otur\'e ou trait\'e.
  \item D\'el\'egation reu ou r\'evoqu\'ee.
\end{itemize}

\begin{steps}
  \item Cliquer sur la cloche \faBell\ pour afficher la liste des notifications.
  \item Cliquer sur une notification pour la marquer comme lue.
  \item Naviguer vers la fonctionnalit\'e concern\'ee via le lien dans la notification.
\end{steps}

\screenshot{img_notifications}{Panneau de notifications}{Capturer le panneau de notifications ouvert, montrant plusieurs notifications (certaines lues, certaines non lues). La cloche avec un badge num\'erique doit \^etre visible.}

\chapter{R\^oles et Droits}
\label{chap:roles}

\section{Pr\'esentation des r\^oles}

L'application CGP met en \oe uvre un contr\^ole d'acc\`es bas\'e sur les r\^oles (\textit{Role-Based Access Control}). Chaque utilisateur se voit attribuer un ou plusieurs r\^oles via des \textbf{affectations}. Un r\^ole est toujours li\'e \`a une \textbf{structure} et \`a une \textbf{ann\'ee acad\'emique}.

\medskip
\begin{notebox}
  Un m\^eme utilisateur peut poss\'eder plusieurs r\^oles dans des structures diff\'erentes, voire dans la m\^eme structure. Les droits dont il dispose correspondent \`a l'union de tous ses r\^oles actifs.
\end{notebox}

\section{Catalogue des r\^oles}

Le tableau suivant pr\'esente l'ensemble des r\^oles disponibles dans l'application, class\'es par niveau d'acc\`es d\'ecroissant.

\medskip
\begin{longtable}{L{4.5cm}L{3cm}X}
  \toprule
  \rowcolor{cgplightblue}
  \textbf{Identifiant technique} & \textbf{Niveau} & \textbf{Description} \\
  \midrule
  \endfirsthead
  \toprule
  \rowcolor{cgplightblue}
  \textbf{Identifiant technique} & \textbf{Niveau} & \textbf{Description} \\
  \midrule
  \endhead
  \midrule
  \multicolumn{3}{r}{\small\textit{(suite page suivante)}}\\
  \endfoot
  \bottomrule
  \endlastfoot

  \texttt{services-centraux} & \cellcolor{cgplightred}Maximal
    & Acc\`es complet sans restriction. G\`ere les ann\'ees, les structures, les imports/exports, l'audit, et tous les param\'etrages. \\
  \rowcolor{cgplightgray}
  \texttt{administrateur} & Administratif
    & Administration du syst\`eme\,: gestion des utilisateurs, imports/exports, d\'el\'egations, organigrammes. \\
  \texttt{directeur-composante} & Direction composante
    & Gestion des utilisateurs et affectations dans sa composante. Import/export, organigrammes, d\'el\'egations dans son p\'erim\`etre. \\
  \rowcolor{cgplightgray}
  \texttt{directeur-administratif} & Direction admin.
    & Gestion administrative\,: utilisateurs, affectations, organigrammes, d\'el\'egations dans son p\'erim\`etre. \\
  \texttt{directeur-administratif-adjoint} & Direction admin. adj.
    & Similaire \`a \texttt{directeur-administratif} avec des droits l\'eg\`erement r\'eduits. \\
  \rowcolor{cgplightgray}
  \texttt{directeur-departement} & Direction d\'ept.
    & Gestion des affectations et d\'el\'egations dans son d\'epartement. G\'en\`ere des organigrammes. \\
  \texttt{vice-president-departement} & VP d\'ept.
    & Similaire au directeur de d\'epartement, avec acc\`es aux organigrammes et d\'el\'egations. \\
  \rowcolor{cgplightgray}
  \texttt{directeur-adjoint-licence} & Direction adj. licence
    & Gestion du p\'erim\`etre licence. Affectations et organigrammes. \\
  \texttt{responsable-service-pedagogique} & Resp. p\'edagogique
    & Gestion des affectations dans le service p\'edagogique. Organigrammes. \\
  \rowcolor{cgplightgray}
  \texttt{responsable-adjoint-service-pedagogique} & Resp. adj. p\'edagogique
    & Droits r\'eduits par rapport au responsable du service p\'edagogique. \\
  \texttt{directeur-mention} & Direction mention
    & Gestion de la mention. Affectations, organigrammes, d\'el\'egations. \\
  \rowcolor{cgplightgray}
  \texttt{directeur-specialite} & Direction sp\'ecialit\'e
    & Gestion de la sp\'ecialit\'e (mention). Affectations, organigrammes. \\
  \texttt{responsable-formation} & Resp. formation
    & Responsable d'une formation. Affectations li\'ees \`a sa formation, organigrammes. \\
  \rowcolor{cgplightgray}
  \texttt{responsable-annee} & Resp. ann\'ee
    & Responsable d'une ann\'ee d'\'etude. Affectations li\'ees, organigrammes. \\
  \texttt{directeur-etudes} & Directeur \'etudes
    & G\`ere les affectations li\'ees aux \'etudes. Organigrammes. \\
  \rowcolor{cgplightgray}
  \texttt{responsable-qualite} & Resp. qualit\'e
    & Acc\`es \`a l'annuaire et aux organigrammes dans son p\'erim\`etre. \\
  \texttt{responsable-international} & Resp. international
    & Acc\`es \`a l'annuaire et aux organigrammes li\'es aux programmes internationaux. \\
  \rowcolor{cgplightgray}
  \texttt{referent-commun} & R\'ef\'erent commun
    & Acc\`es g\'en\'eral \`a l'annuaire et aux organigrammes. \\
  \texttt{directeur-adjoint-ecole} & Directeur adj. \'ecole
    & Gestion du p\'erim\`etre \'ecole. Affectations et organigrammes. \\
  \rowcolor{cgplightgray}
  \texttt{secretariat-pedagogique} & Secr\'etariat
    & Consultation de l'annuaire. Cr\'eation et mise \`a jour des contacts de r\^ole. Signalements. \\
  \texttt{utilisateur-simple} & Utilisateur
    & Consultation de l'annuaire et des organigrammes. Signalements. \\
  \rowcolor{cgplightgray}
  \texttt{lecture-seule} & Lecture seule
    & Consultation uniquement. Aucune modification. Signalements. \\

\end{longtable}

\section{Regroupement fonctionnel des r\^oles}

Pour faciliter la lecture de ce manuel, les r\^oles sont regroup\'es en \textbf{profils fonctionnels} correspondant aux chapitres d\'edi\'es\,:

\medskip
\begin{tabularx}{\textwidth}{L{3.5cm}XC{2.5cm}}
  \toprule
  \rowcolor{cgplightblue}
  \textbf{Profil fonctionnel} & \textbf{R\^oles inclus} & \textbf{Chapitre} \\
  \midrule
  Services Centraux
    & \texttt{services-centraux}
    & Chap.~\ref{chap:sc} \\
  \rowcolor{cgplightgray}
  Direction composante / admin.
    & \texttt{directeur-composante}, \texttt{administrateur}, \texttt{directeur-administratif}, \texttt{directeur-administratif-adjoint}
    & Chap.~\ref{chap:dircomp} \\
  Directions et responsables
    & \texttt{directeur-departement}, \texttt{vice-president-departement}, \texttt{directeur-adjoint-licence}, \texttt{responsable-service-pedagogique}, \texttt{responsable-adjoint-service-pedagogique}, \texttt{directeur-mention}, \texttt{directeur-specialite}, \texttt{responsable-formation}, \texttt{responsable-annee}, \texttt{directeur-etudes}, \texttt{responsable-qualite}, \texttt{responsable-international}, \texttt{referent-commun}, \texttt{directeur-adjoint-ecole}
    & Chap.~\ref{chap:directions} \\
  \rowcolor{cgplightgray}
  Secr\'etariat p\'edagogique
    & \texttt{secretariat-pedagogique}
    & Chap.~\ref{chap:secretariat} \\
  Utilisateur standard
    & \texttt{utilisateur-simple}, \texttt{lecture-seule}
    & Chap.~\ref{chap:utilisateur} \\
  \bottomrule
\end{tabularx}

\section{Matrice des droits}

Le tableau suivant synth\'etise les droits par fonctionnalit\'e et par profil. \\
L\'egende\,: \textbf{\textcolor{cgpgreen}{\checkmark}} = autoris\'e, \textbf{\textcolor{cgpred}{$\times$}} = refus\'e, \textbf{\textcolor{cgporange}{\textasciitilde}} = partiel (p\'erim\`etre limit\'e).

\medskip
\begin{longtable}{L{4.5cm}C{1.6cm}C{1.9cm}C{2cm}C{1.8cm}C{1.5cm}}
  \toprule
  \rowcolor{cgplightblue}
  \textbf{Fonctionnalit\'e}
    & \textbf{SC}
    & \textbf{Dir. comp.}
    & \textbf{Directions}
    & \textbf{Secr\'et.}
    & \textbf{Util.} \\
  \midrule
  \endfirsthead
  \toprule
  \rowcolor{cgplightblue}
  \textbf{Fonctionnalit\'e}
    & \textbf{SC}
    & \textbf{Dir. comp.}
    & \textbf{Directions}
    & \textbf{Secr\'et.}
    & \textbf{Util.} \\
  \midrule
  \endhead
  \bottomrule
  \endlastfoot

  Consulter l'annuaire
    & {\color{cgpgreen}\checkmark}
    & {\color{cgpgreen}\checkmark}
    & {\color{cgpgreen}\checkmark}
    & {\color{cgpgreen}\checkmark}
    & {\color{cgpgreen}\checkmark} \\
  \rowcolor{cgplightgray}
  Voir les organigrammes
    & {\color{cgpgreen}\checkmark}
    & {\color{cgpgreen}\checkmark}
    & {\color{cgpgreen}\checkmark}
    & {\color{cgpgreen}\checkmark}
    & {\color{cgpgreen}\checkmark} \\
  G\'en\'erer des organigrammes
    & {\color{cgpgreen}\checkmark}
    & {\color{cgporange}\textasciitilde}
    & {\color{cgporange}\textasciitilde}
    & {\color{cgpred}$\times$}
    & {\color{cgpred}$\times$} \\
  \rowcolor{cgplightgray}
  Geler/d\'egeler un organigramme
    & {\color{cgpgreen}\checkmark}
    & {\color{cgpred}$\times$}
    & {\color{cgpred}$\times$}
    & {\color{cgpred}$\times$}
    & {\color{cgpred}$\times$} \\
  Cr\'eer des utilisateurs
    & {\color{cgpgreen}\checkmark}
    & {\color{cgporange}\textasciitilde}
    & {\color{cgpred}$\times$}
    & {\color{cgpred}$\times$}
    & {\color{cgpred}$\times$} \\
  \rowcolor{cgplightgray}
  Modifier des utilisateurs
    & {\color{cgpgreen}\checkmark}
    & {\color{cgporange}\textasciitilde}
    & {\color{cgpred}$\times$}
    & {\color{cgpred}$\times$}
    & {\color{cgpred}$\times$} \\
  Supprimer des utilisateurs
    & {\color{cgpgreen}\checkmark}
    & {\color{cgporange}\textasciitilde}
    & {\color{cgpred}$\times$}
    & {\color{cgpred}$\times$}
    & {\color{cgpred}$\times$} \\
  \rowcolor{cgplightgray}
  G\'erer les affectations
    & {\color{cgpgreen}\checkmark}
    & {\color{cgporange}\textasciitilde}
    & {\color{cgporange}\textasciitilde}
    & {\color{cgpred}$\times$}
    & {\color{cgpred}$\times$} \\
  G\'erer les contacts de r\^ole
    & {\color{cgpgreen}\checkmark}
    & {\color{cgpgreen}\checkmark}
    & {\color{cgpgreen}\checkmark}
    & {\color{cgpgreen}\checkmark}
    & {\color{cgpred}$\times$} \\
  \rowcolor{cgplightgray}
  Modifier les structures
    & {\color{cgpgreen}\checkmark}
    & {\color{cgpred}$\times$}
    & {\color{cgpred}$\times$}
    & {\color{cgpred}$\times$}
    & {\color{cgpred}$\times$} \\
  Importer des donn\'ees
    & {\color{cgpgreen}\checkmark}
    & {\color{cgporange}\textasciitilde}
    & {\color{cgpred}$\times$}
    & {\color{cgpred}$\times$}
    & {\color{cgpred}$\times$} \\
  \rowcolor{cgplightgray}
  Exporter des donn\'ees
    & {\color{cgpgreen}\checkmark}
    & {\color{cgporange}\textasciitilde}
    & {\color{cgpred}$\times$}
    & {\color{cgpred}$\times$}
    & {\color{cgpred}$\times$} \\
  G\'erer les ann\'ees acad\'emiques
    & {\color{cgpgreen}\checkmark}
    & {\color{cgpred}$\times$}
    & {\color{cgpred}$\times$}
    & {\color{cgpred}$\times$}
    & {\color{cgpred}$\times$} \\
  \rowcolor{cgplightgray}
  Changer l'ann\'ee active
    & {\color{cgpgreen}\checkmark}
    & {\color{cgpred}$\times$}
    & {\color{cgpred}$\times$}
    & {\color{cgpred}$\times$}
    & {\color{cgpred}$\times$} \\
  G\'erer les d\'el\'egations
    & {\color{cgpgreen}\checkmark}
    & {\color{cgpgreen}\checkmark}
    & {\color{cgpgreen}\checkmark}
    & {\color{cgpred}$\times$}
    & {\color{cgpred}$\times$} \\
  \rowcolor{cgplightgray}
  Soumettre demande de r\^ole
    & {\color{cgpred}$\times$}
    & {\color{cgpgreen}\checkmark}
    & {\color{cgpgreen}\checkmark}
    & {\color{cgpred}$\times$}
    & {\color{cgpred}$\times$} \\
  Traiter demandes de r\^oles
    & {\color{cgpgreen}\checkmark}
    & {\color{cgpred}$\times$}
    & {\color{cgpred}$\times$}
    & {\color{cgpred}$\times$}
    & {\color{cgpred}$\times$} \\
  \rowcolor{cgplightgray}
  Cr\'eer un signalement
    & {\color{cgpgreen}\checkmark}
    & {\color{cgpgreen}\checkmark}
    & {\color{cgpgreen}\checkmark}
    & {\color{cgpgreen}\checkmark}
    & {\color{cgpgreen}\checkmark} \\
  Traiter les signalements
    & {\color{cgpgreen}\checkmark}
    & {\color{cgporange}\textasciitilde}
    & {\color{cgporange}\textasciitilde}
    & {\color{cgpred}$\times$}
    & {\color{cgpred}$\times$} \\
  \rowcolor{cgplightgray}
  Consulter le journal d'audit
    & {\color{cgpgreen}\checkmark}
    & {\color{cgpred}$\times$}
    & {\color{cgpred}$\times$}
    & {\color{cgpred}$\times$}
    & {\color{cgpred}$\times$} \\
  Modifier son propre profil
    & {\color{cgpgreen}\checkmark}
    & {\color{cgpgreen}\checkmark}
    & {\color{cgpgreen}\checkmark}
    & {\color{cgpgreen}\checkmark}
    & {\color{cgpgreen}\checkmark} \\

\end{longtable}

\medskip
\begin{notebox}[Droits partiels (\textasciitilde)]
  La mention \og partiel \fg{} signifie que le droit existe mais est limit\'e au \textbf{p\'erim\`etre de la structure} rattach\'ee \`a l'affectation de l'utilisateur. Par exemple, un directeur de composante ne peut g\'erer les utilisateurs que dans sa composante et ses sous-structures.
\end{notebox}

\section{P\'erim\`etre et hi\'erarchie des structures}

La hi\'erarchie des structures acad\'emiques dans CGP est la suivante (du plus g\'en\'eral au plus sp\'ecifique)\,:

\medskip
\begin{center}
\begin{tikzpicture}[
  node distance=0.5cm,
  box/.style={rectangle, rounded corners=4pt, draw=cgpblue, fill=cgplightblue,
              text width=3.5cm, align=center, font=\small\bfseries, minimum height=0.8cm}
]
  \node[box] (comp) {Composante};
  \node[box, below=of comp] (dept) {D\'epartement};
  \node[box, below=of dept] (mention) {Mention};
  \node[box, below=of mention] (parcours) {Parcours};
  \node[box, below=of parcours] (niveau) {Niveau};
  \draw[->, thick, cgpblue] (comp) -- (dept);
  \draw[->, thick, cgpblue] (dept) -- (mention);
  \draw[->, thick, cgpblue] (mention) -- (parcours);
  \draw[->, thick, cgpblue] (parcours) -- (niveau);
\end{tikzpicture}
\end{center}

\medskip
Un utilisateur ayant un r\^ole dans une structure donn\'ee acc\`ede g\'en\'eralement \`a cette structure \textbf{et \`a l'ensemble de ses descendants} dans la hi\'erarchie. Les Services Centraux n'ont pas de restriction de p\'erim\`etre.

\begin{tipbox}[P\'erim\`etre de consultation versus de g\'en\'eration]
  Pour les organigrammes, la \textbf{consultation} est plus large que la \textbf{g\'en\'eration}\,: un utilisateur peut consulter des organigrammes d\'ej\`a g\'en\'er\'es pour des structures hors de son p\'erim\`etre, mais il ne peut en g\'en\'erer que pour son p\'erim\`etre propre.
\end{tipbox}

\chapter{Manuel --- Services Centraux}
\label{chap:sc}

\begin{rolebox}{services-centraux}
  Ce chapitre d\'ecrit \textbf{l'ensemble} des fonctionnalit\'es accessibles au r\^ole \texttt{services-centraux}, qui dispose d'un acc\`es complet et non restreint \`a l'application.
\end{rolebox}

\section{Tableau de bord}

Le tableau de bord est la page d'accueil apr\`es connexion. Pour les Services Centraux, il affiche\,:

\begin{itemize}
  \item Le \textbf{s\'electeur d'ann\'ee acad\'emique} permettant de basculer entre toutes les ann\'ees disponibles.
  \item Des \textbf{statistiques globales} (nombre de structures, utilisateurs, affectations, signalements en attente).
  \item Des \textbf{raccourcis} vers toutes les rubriques du menu.
  \item Les \textbf{notifications} non lues.
\end{itemize}

\screenshot{img_dashboard_sc}{Tableau de bord -- Services Centraux}{Capturer la page tableau de bord avec un compte services-centraux. Montrer\,: le s\'electeur d'ann\'ee ouvert (menu d\'eroulant visible), les statistiques chiffr\'ees, les raccourcis de navigation et la cloche de notification avec un badge.}

\section{Gestion des ann\'ees acad\'emiques}

La rubrique \textbf{Ann\'ees} (accessible via le menu) permet de g\'erer le cycle de vie des ann\'ees acad\'emiques.

\screenshot{img_annees_liste}{Liste des ann\'ees acad\'emiques}{Capturer la page de gestion des ann\'ees. Montrer la liste des ann\'ees existantes avec leurs statuts (EN\_COURS, PREPARATION, ARCHIVEE) repr\'esent\'es par des badges de couleur, et les boutons d'action disponibles pour chaque ann\'ee.}

\subsection{Statuts des ann\'ees}

Chaque ann\'ee poss\`ede un statut\,:

\medskip
\begin{tabularx}{\textwidth}{L{3cm}C{2.5cm}X}
  \toprule
  \rowcolor{cgplightblue}
  \textbf{Statut} & \textbf{Couleur} & \textbf{Signification} \\
  \midrule
  \texttt{EN\_COURS}   & \cellcolor{cgplightgreen}Vert  & Ann\'ee active visible par tous les utilisateurs. Une seule ann\'ee peut avoir ce statut. \\
  \rowcolor{cgplightgray}
  \texttt{PREPARATION} & \cellcolor{cgplightorange}Orange & Ann\'ee en pr\'eparation, non encore active. Visible uniquement pour les SC. \\
  \texttt{ARCHIVEE}    & \cellcolor{cgplightgray}Gris  & Ann\'ee cl\^otur\'ee. Donn\'ees consultables mais non modifiables (sauf par SC). \\
  \bottomrule
\end{tabularx}

\subsection{Cr\'eer une nouvelle ann\'ee}

\begin{steps}
  \item Naviguer vers \textbf{Ann\'ees} dans le menu.
  \item Cliquer sur le bouton \textbf{Cr\'eer une ann\'ee}.
  \item Remplir le formulaire\,:
    \begin{itemize}
      \item \textit{Libell\'e} : intitul\'e de l'ann\'ee (ex. : \og 2025-2026 \fg).
      \item \textit{Statut initial} : PREPARATION (recommand\'e) ou EN\_COURS.
    \end{itemize}
  \item Cliquer sur \textbf{Cr\'eer}.
\end{steps}

\screenshot{img_annee_creer}{Formulaire de cr\'eation d'ann\'ee}{Capturer le formulaire de cr\'eation d'une nouvelle ann\'ee acad\'emique avec tous les champs visibles. Le formulaire peut \^etre une modale ou un panneau lat\'eral selon l'interface.}

\subsection{Cloner une ann\'ee existante}

Le clonage permet de cr\'eer une nouvelle ann\'ee en copiant tout ou partie des donn\'ees d'une ann\'ee source (structures, affectations).

\begin{steps}
  \item Sur la liste des ann\'ees, cliquer sur \textbf{Cloner} (ou l'ic\^one correspondante) \`a c\^ot\'e de l'ann\'ee source.
  \item Configurer le clonage\,:
    \begin{itemize}
      \item \textit{Libell\'e} de la nouvelle ann\'ee.
      \item \textit{Structures \`a cloner} : toutes ou s\'election partielle.
      \item \textit{Copier les affectations} : oui ou non (si non, seules les structures sont copi\'ees).
    \end{itemize}
  \item Cliquer sur \textbf{Confirmer le clonage}.
  \item L'application cr\'ee la nouvelle ann\'ee et affiche une confirmation.
\end{steps}

\screenshot{img_annee_cloner}{Formulaire de clonage d'ann\'ee}{Capturer le formulaire de clonage d'ann\'ee, montrant\,: le champ libell\'e, la liste de s\'election des structures \`a inclure (cases \`a cocher), et l'option de copie des affectations. Si c'est une modale, capturer la modale enti\`ere.}

\begin{tipbox}[Cloner sans les affectations]
  Le clonage sans les affectations est utile en d\'ebut d'ann\'ee lorsque les r\'eorganisations sont fr\'equentes. Les structures seront en place, mais les responsables devront \^etre affect\'es manuellement.
\end{tipbox}

\subsection{Changer le statut d'une ann\'ee}

\begin{steps}
  \item Sur la liste des ann\'ees, localiser l'ann\'ee souhait\'ee.
  \item Cliquer sur le bouton \textbf{Activer} (pour passer en EN\_COURS) ou \textbf{Archiver} selon la transition souhait\'ee.
  \item Confirmer le changement dans la bo\^ite de dialogue.
\end{steps}

\begin{warnbox}[Passage en EN\_COURS]
  Activer une ann\'ee en EN\_COURS archive automatiquement l'ann\'ee pr\'ec\'edemment active. Cette action affecte tous les utilisateurs\,: ils basculeront sur la nouvelle ann\'ee courante.
\end{warnbox}

\subsection{Supprimer une ann\'ee}

\begin{steps}
  \item Sur la liste des ann\'ees, cliquer sur \textbf{Supprimer} pour l'ann\'ee concern\'ee.
  \item L'application g\'en\`ere \textbf{automatiquement une sauvegarde} (export Excel) avant suppression.
  \item T\'el\'echarger la sauvegarde propos\'ee.
  \item Confirmer la suppression en saisissant le libell\'e de l'ann\'ee si demand\'e.
\end{steps}

\begin{warnbox}[Suppression irr\'eversible]
  La suppression d'une ann\'ee est \textbf{d\'efinitive}. Assurez-vous d'avoir t\'el\'echarg\'e et archiv\'e la sauvegarde avant de confirmer.
\end{warnbox}

%% ============================================================
\section{Gestion des utilisateurs}
%% ============================================================

La rubrique \textbf{Responsables} (ou via la gestion des utilisateurs dans les param\`etres) permet de cr\'eer, modifier et supprimer des comptes utilisateurs.

\screenshot{img_users_liste}{Liste des utilisateurs}{Capturer la page de gestion des utilisateurs/responsables avec la barre de recherche, les filtres hi\'erarchiques (composante, d\'epartement, etc.) et le tableau ou la liste de r\'esultats. Montrer au moins 5 utilisateurs diff\'erents.}

\subsection{Cr\'eer un utilisateur}

\begin{steps}
  \item Naviguer vers \textbf{Responsables} dans le menu.
  \item Cliquer sur \textbf{Cr\'eer un utilisateur}.
  \item Remplir le formulaire\,:
    \begin{itemize}
      \item \textit{Nom} (obligatoire)
      \item \textit{Pr\'enom} (obligatoire)
      \item \textit{Identifiant / Login} (obligatoire, unique dans le syst\`eme)
      \item \textit{Courriel institutionnel}
      \item \textit{T\'el\'ephone}
      \item \textit{Bureau}
    \end{itemize}
  \item Cliquer sur \textbf{Enregistrer}.
  \item L'utilisateur est cr\'e\'e. Il faut ensuite lui attribuer une ou plusieurs affectations.
\end{steps}

\screenshot{img_user_creer}{Formulaire de cr\'eation d'utilisateur}{Capturer le formulaire de cr\'eation d'un utilisateur avec tous les champs visibles. Le formulaire peut \^etre une page d\'edi\'ee, une modale ou un panneau lat\'eral.}

\subsection{Modifier un utilisateur}

\begin{steps}
  \item Rechercher l'utilisateur via la barre de recherche (nom, pr\'enom, identifiant).
  \item Cliquer sur l'utilisateur dans la liste des r\'esultats.
  \item Cliquer sur \textbf{Modifier}.
  \item Mettre \`a jour les champs souhait\'es.
  \item Cliquer sur \textbf{Enregistrer}.
\end{steps}

\begin{notebox}
  Un utilisateur ne peut pas modifier lui-m\^eme son propre nom, pr\'enom ou identifiant. Ces champs sont r\'eserv\'es aux Services Centraux et aux directeurs de composante.
\end{notebox}

\subsection{Supprimer un utilisateur}

\begin{steps}
  \item Rechercher l'utilisateur.
  \item Cliquer sur l'utilisateur puis sur \textbf{Supprimer}.
  \item Confirmer la suppression.
\end{steps}

\begin{warnbox}
  La suppression d'un utilisateur entra\^ine la suppression de toutes ses affectations. Cette action est \textbf{irr\'eversible}.
\end{warnbox}

%% ============================================================
\section{Gestion des structures}
%% ============================================================

\begin{rolebox}{services-centraux}
  Seuls les Services Centraux peuvent \textbf{modifier} les informations d'une structure. Tous les utilisateurs peuvent les consulter.
\end{rolebox}

\screenshot{img_structures_liste}{Liste des structures}{Capturer la page de gestion des structures avec la barre de recherche et le tableau listant les structures (nom, type, parent, ann\'ee). Montrer plusieurs types de structures diff\'erentes.}

\subsection{Consulter une fiche structure}

\begin{steps}
  \item Naviguer vers \textbf{Structures} dans le menu.
  \item Rechercher la structure par son nom ou son code.
  \item Cliquer sur la structure pour ouvrir sa fiche d\'etaill\'ee.
\end{steps}

La fiche structure affiche\,:
\begin{itemize}
  \item \textbf{Identification} : nom, type, code, ann\'ee.
  \item \textbf{Hi\'erarchie} : structure parente et sous-structures directes.
  \item \textbf{Coordonn\'ees} : adresse, t\'el\'ephone, courriel fonctionnel.
  \item \textbf{Responsables} : liste des personnes affect\'ees avec leurs r\^oles.
  \item \textbf{Secr\'etariat} : contacts du secr\'etariat p\'edagogique.
  \item \textbf{Statistiques} : nombre d'affectations, de sous-structures, etc.
\end{itemize}

\screenshot{img_structure_detail}{Fiche d\'etaill\'ee d'une structure}{Capturer la fiche compl\`ete d'une structure (composante ou d\'epartement). Montrer toutes les sections\,: identification, hi\'erarchie, responsables, secr\'etariat. Faire d\'efiler si n\'ecessaire et capturer l'ensemble ou faire plusieurs captures pour chaque section.}

\subsection{Modifier une structure}

\begin{steps}
  \item Ouvrir la fiche de la structure \`a modifier.
  \item Cliquer sur \textbf{Modifier} (ou l'ic\^one crayon).
  \item Mettre \`a jour les champs autoris\'es\,: nom, coordonn\'ees, etc.
  \item Cliquer sur \textbf{Enregistrer}.
\end{steps}

\screenshot{img_structure_modifier}{Formulaire de modification d'une structure}{Capturer le formulaire d'\'edition d'une structure en mode modification. Tous les champs \'editables doivent \^etre visibles (nom, coordonn\'ees, adresse, courriel fonctionnel).}

%% ============================================================
\section{Gestion des affectations}
%% ============================================================

Les affectations lient un \textbf{utilisateur} \`a un \textbf{r\^ole} dans une \textbf{structure} pour une \textbf{ann\'ee} donn\'ee.

\screenshot{img_affectations_liste}{Interface de gestion des affectations}{Capturer la page de gestion des responsables/affectations avec les filtres hi\'erarchiques (composante, d\'epartement, mention, parcours, niveau) et les r\'esultats de recherche. Montrer une liste d'affectations avec nom, r\^ole et structure.}

\subsection{Cr\'eer une affectation}

\begin{steps}
  \item Naviguer vers \textbf{Responsables} dans le menu.
  \item Rechercher l'utilisateur concern\'e.
  \item Dans la fiche de l'utilisateur, cliquer sur \textbf{Ajouter un r\^ole} (ou \textbf{Nouvelle affectation}).
  \item S\'electionner\,:
    \begin{itemize}
      \item \textit{R\^ole} dans la liste des r\^oles disponibles.
      \item \textit{Composante} (obligatoire).
      \item \textit{D\'epartement}, \textit{Mention}, \textit{Parcours}, \textit{Niveau} (selon le r\^ole et la granularit\'e souhait\'ee).
      \item \textit{N+1} : le superviseur hi\'erarchique direct (optionnel).
    \end{itemize}
  \item Cliquer sur \textbf{Enregistrer}.
\end{steps}

\screenshot{img_affectation_creer}{Formulaire de cr\'eation d'affectation}{Capturer le formulaire de cr\'eation d'affectation avec tous les menus d\'eroulants visibles\,: r\^ole, composante, d\'epartement, mention, parcours, niveau, et superviseur N+1. Si certains champs sont conditionnels, montrer un \'etat avec plusieurs niveaux s\'electionn\'es.}

\subsection{Modifier ou supprimer une affectation}

\begin{steps}
  \item Localiser l'affectation dans la liste (rechercher par nom, r\^ole ou structure).
  \item Cliquer sur l'affectation.
  \item Pour modifier : cliquer sur \textbf{Modifier}, apporter les changements, puis \textbf{Enregistrer}.
  \item Pour supprimer : cliquer sur \textbf{Supprimer} et confirmer.
\end{steps}

\subsection{G\'erer les contacts de r\^ole}

Un contact de r\^ole est une information de contact fonctionnelle associ\'ee \`a une affectation\,: courriel fonctionnel, t\'el\'ephone de service, bureau.

\begin{steps}
  \item Localiser l'affectation concern\'ee.
  \item Cliquer sur \textbf{Contact de r\^ole} (ou \textbf{\'Editer le contact}).
  \item Renseigner\,:
    \begin{itemize}
      \item \textit{Courriel fonctionnel} (ex. : \texttt{dir.composante@univ.fr})
      \item \textit{T\'el\'ephone de service}
      \item \textit{Bureau fonctionnel}
    \end{itemize}
  \item Cliquer sur \textbf{Enregistrer}.
\end{steps}

\screenshot{img_affectation_contact}{Formulaire de contact de r\^ole}{Capturer le formulaire d'\'edition du contact de r\^ole d'une affectation (courriel fonctionnel, t\'el\'ephone, bureau). Ce formulaire appara\^it g\'en\'eralement dans un panneau lat\'eral ou une modale.}

%% ============================================================
\section{Import et Export de donn\'ees}
%% ============================================================

La rubrique \textbf{Import / Export} permet d'\'echanger des donn\'ees via des classeurs Excel standardis\'es (format \texttt{CGP\_STANDARD\_V1}).

\screenshot{img_import_export}{Page Import / Export}{Capturer la page Import/Export enti\`ere. Montrer les deux zones principales\,: la zone Export (avec les options de s\'election d'ann\'ee et de structure) et la zone Import (avec la zone de d\'ep\^ot de fichier).}

\subsection{Exporter un classeur}

\begin{steps}
  \item Naviguer vers \textbf{Import / Export}.
  \item Dans la section \textbf{Export}\,:
    \begin{itemize}
      \item S\'electionner l'\textit{Ann\'ee} \`a exporter.
      \item Optionnel : s\'electionner une \textit{Structure} sp\'ecifique (laisser vide pour exporter toute l'ann\'ee).
      \item Choisir le \textit{Format} (Excel recommand\'e).
    \end{itemize}
  \item Cliquer sur \textbf{T\'el\'echarger}.
  \item Le fichier Excel est t\'el\'echarg\'e dans le dossier de t\'el\'echargements du navigateur.
\end{steps}

\subsection{Exporter un mod\`ele vierge}

Pour obtenir un mod\`ele de classeur vide (pour saisie manuelle) :

\begin{steps}
  \item Dans la section Export, cliquer sur \textbf{T\'el\'echarger le mod\`ele vierge}.
  \item Remplir le mod\`ele selon les instructions contenues dans les onglets du classeur.
\end{steps}

\subsection{Importer un classeur}

\begin{steps}
  \item Dans la section \textbf{Import}, cliquer sur \textbf{S\'electionner un fichier} ou faire glisser le fichier \texttt{.xlsx} dans la zone de d\'ep\^ot.
  \item Optionnel : s\'electionner une \textit{Structure cible} pour limiter l'import \`a un p\'erim\`etre.
  \item Cliquer sur \textbf{Pr\'evisualiser}.
  \item L'application analyse le fichier et affiche un r\'esum\'e des op\'erations pr\'evues\,:
    \begin{itemize}
      \item \textbf{Cr\'eations} : nouvelles entr\'ees \`a ins\'erer.
      \item \textbf{Mises \`a jour} : entr\'ees existantes modifi\'ees.
      \item \textbf{Conflits} : entr\'ees n\'ecessitant une d\'ecision.
    \end{itemize}
  \item Examiner la pr\'evisualisation. R\'esoudre les conflits si n\'ecessaire.
  \item Cliquer sur \textbf{Confirmer l'import}.
\end{steps}

\screenshot{img_import_preview}{Pr\'evisualisation d'import}{Capturer la zone de pr\'evisualisation d'import affich\'ee apr\`es analyse du fichier. Montrer un tableau avec des lignes colori\'ees selon leur statut (vert = cr\'eation, orange = mise \`a jour, rouge = conflit). Montrer au moins un exemple de chaque statut si possible.}

\begin{warnbox}[V\'erifier avant de confirmer]
  Une fois l'import confirm\'e, les donn\'ees sont \'ecrites en base. Lisez attentivement la pr\'evisualisation, en particulier les conflits, avant de cliquer sur \textbf{Confirmer}.
\end{warnbox}

\begin{tipbox}[Importation par p\'erim\`etre]
  Pour ne mettre \`a jour qu'une composante sp\'ecifique, utilisez le filtre \og Structure cible \fg{}. Les donn\'ees hors de ce p\'erim\`etre ne seront pas modifi\'ees.
\end{tipbox}

%% ============================================================
\section{Organigrammes}
%% ============================================================

La rubrique \textbf{Organigramme} permet de g\'en\'erer, consulter, geler et exporter des organigrammes hi\'erarchiques.

\subsection{G\'en\'erer un organigramme}

\begin{steps}
  \item Naviguer vers \textbf{Organigramme}.
  \item S\'electionner le \textbf{type de vue}\,:
    \begin{itemize}
      \item \textbf{Vue Structures} : repr\'esentation hi\'erarchique des entit\'es (composantes, d\'epartements, mentions...).
      \item \textbf{Vue Personnes} : repr\'esentation hi\'erarchique des personnes (N+1).
    \end{itemize}
  \item Choisir la \textbf{structure racine} (point de d\'epart de l'organigramme).
  \item Appliquer les filtres souhaait\'es (r\^ole, composante, d\'epartement, mention, parcours, niveau).
  \item Cliquer sur \textbf{G\'en\'erer}.
  \item L'organigramme s'affiche. S'il existait d\'ej\`a un organigramme identique non gel\'e, il est r\'eutilis\'e.
\end{steps}

\screenshot{img_organigramme_structures}{Vue Structures de l'organigramme}{Capturer l'organigramme en vue Structures avec plusieurs niveaux de hi\'erarchie visibles (composante, d\'epartements, mentions). Le panneau de contr\^ole avec les filtres doit \^etre visible.}

\screenshot{img_organigramme_personnes}{Vue Personnes de l'organigramme}{Capturer l'organigramme en vue Personnes avec des n\oe uds repr\'esentant des personnes. Les filtres (r\^ole, composante, etc.) doivent \^etre visibles. Montrer un arbre avec au moins 3 niveaux hi\'erarchiques N+1.}

\subsection{Geler et d\'egeler un organigramme}

\begin{rolebox}{services-centraux}
  Le gel et le d\'egel sont r\'eserv\'es aux Services Centraux.
\end{rolebox}

Geler un organigramme le fige dans son \'etat actuel. Il ne sera plus remplac\'e lors des prochaines g\'en\'erations.

\begin{steps}
  \item Consulter l'organigramme concern\'e.
  \item Cliquer sur le bouton \textbf{Geler} dans la barre d'actions.
  \item Confirmer. L'organigramme affiche le badge \og Gel\'e \fg.
\end{steps}

Pour d\'egeler\,:
\begin{steps}
  \item Ouvrir l'organigramme gel\'e.
  \item Cliquer sur \textbf{D\'egeler}.
  \item L'organigramme redevient disponible pour remplacement automatique.
\end{steps}

\subsection{Biblioth\`eque des organigrammes}

La biblioth\`eque liste tous les organigrammes g\'en\'er\'es, avec des filtres par structure, type de vue et statut (gel\'e / disponible).

\begin{steps}
  \item Dans la page \textbf{Organigramme}, faire d\'efiler vers la section \og Biblioth\`eque \fg{} ou cliquer sur l'onglet correspondant.
  \item Utiliser les filtres pour retrouver l'organigramme souhait\'e.
  \item Cliquer sur \textbf{Voir en structures} ou \textbf{Voir en personnes} pour l'afficher.
\end{steps}

\screenshot{img_organigramme_biblio}{Biblioth\`eque des organigrammes g\'en\'er\'es}{Capturer la section biblioth\`eque avec les filtres et la liste des organigrammes. Montrer des badges de statut (Gel\'e / Disponible) et les boutons d'action \og Voir en structures \fg{} et \og Voir en personnes \fg{}.}

\subsection{Exporter un organigramme}

\begin{steps}
  \item Afficher l'organigramme souhait\'e.
  \item Cliquer sur le bouton \textbf{Exporter}.
  \item Choisir le format\,: \textbf{PDF}, \textbf{CSV}, \textbf{SVG} ou \textbf{JSON}.
  \item Le fichier est t\'el\'echarg\'e automatiquement.
\end{steps}

\screenshot{img_organigramme_export}{Options d'export de l'organigramme}{Capturer le menu ou la modale d'export de l'organigramme avec les diff\'erents formats disponibles (PDF, CSV, SVG, JSON).}

%% ============================================================
\section{Demandes de r\^oles}
%% ============================================================

La rubrique \textbf{Demandes de r\^oles} permet aux Services Centraux d'examiner et de traiter les demandes soumises par les directions.

\screenshot{img_demandes_roles_sc}{Liste des demandes de r\^oles (Services Centraux)}{Capturer la page de gestion des demandes de r\^oles avec la liste des demandes (en attente, approuv\'ees, refus\'ees). Montrer les informations visibles par les SC\,: demandeur, r\^ole demand\'e, structure, statut, date.}

\subsection{Examiner une demande}

\begin{steps}
  \item Naviguer vers \textbf{Demandes de r\^oles}.
  \item Filtrer par statut \og En attente \fg{} pour voir les demandes \`a traiter.
  \item Cliquer sur une demande pour afficher ses d\'etails.
  \item Lire le motif de la demande et les informations du demandeur.
\end{steps}

\subsection{Approuver ou refuser une demande}

\begin{steps}
  \item Dans le d\'etail de la demande, cliquer sur \textbf{Approuver} ou \textbf{Refuser}.
  \item Si refus, renseigner un \textit{commentaire de justification} (recommand\'e).
  \item Confirmer. Le demandeur re\c{c}oit une notification du r\'esultat.
\end{steps}

\screenshot{img_demande_role_traitement}{Interface de traitement d'une demande de r\^ole}{Capturer le d\'etail d'une demande de r\^ole avec les boutons Approuver/Refuser visibles, le champ commentaire et les informations de la demande (demandeur, r\^ole, structure, motif).}

%% ============================================================
\section{D\'el\'egations}
%% ============================================================

La rubrique \textbf{D\'el\'egations} permet de transf\'erer temporairement des droits d'un utilisateur \`a un autre.

\screenshot{img_delegations_liste}{Liste des d\'el\'egations}{Capturer la page de gestion des d\'el\'egations avec la liste des d\'el\'egations actives (d\'el\'eguant, d\'el\'egu\'e, p\'erim\`etre, dates). Montrer les boutons d'action (Cr\'eer, R\'evoquer).}

\subsection{Cr\'eer une d\'el\'egation}

\begin{steps}
  \item Naviguer vers \textbf{D\'el\'egations}.
  \item Cliquer sur \textbf{Cr\'eer une d\'el\'egation}.
  \item Remplir le formulaire\,:
    \begin{itemize}
      \item \textit{D\'el\'egu\'e} : utilisateur qui re\c{c}oit les droits.
      \item \textit{Structure concern\'ee} : p\'erim\`etre de la d\'el\'egation.
      \item \textit{Date de d\'ebut} et \textit{date de fin}.
    \end{itemize}
  \item Cliquer sur \textbf{Enregistrer}.
\end{steps}

\screenshot{img_delegation_creer}{Formulaire de cr\'eation de d\'el\'egation}{Capturer le formulaire de cr\'eation d'une d\'el\'egation avec les champs d\'el\'egu\'e (recherche d'utilisateur), structure concern\'ee, dates de d\'ebut et de fin.}

\subsection{R\'evoquer une d\'el\'egation}

\begin{steps}
  \item Dans la liste des d\'el\'egations, cliquer sur la d\'el\'egation \`a r\'evoquer.
  \item Cliquer sur \textbf{R\'evoquer}.
  \item Confirmer. La d\'el\'egation est imm\'ediatement annul\'ee.
\end{steps}

\subsection{Exporter la liste des d\'el\'egations}

\begin{steps}
  \item Sur la page D\'el\'egations, cliquer sur \textbf{Exporter (CSV)}.
  \item Le fichier CSV est t\'el\'echarg\'e avec la liste compl\`ete des d\'el\'egations.
\end{steps}

%% ============================================================
\section{Signalements}
%% ============================================================

La rubrique \textbf{Signalements} permet de consulter et de traiter les erreurs signal\'ees par les utilisateurs.

\screenshot{img_signalements_sc}{Liste des signalements (Services Centraux)}{Capturer la page des signalements avec les filtres de statut (En attente, En cours, Cl\^otur\'e), les filtres de type d'erreur (nom, pr\'enom, t\'el\'ephone, bureau, courriel, r\^ole, autre) et la liste des signalements.}

\subsection{Traiter un signalement}

\begin{steps}
  \item Naviguer vers \textbf{Signalements}.
  \item Filtrer par statut \og En attente \fg{}.
  \item Cliquer sur un signalement pour afficher ses d\'etails.
  \item Corriger les donn\'ees dans l'application (aller vers la fiche concern\'ee).
  \item Revenir sur le signalement et cliquer sur \textbf{Cl\^oturer}.
  \item Optionnel : ajouter un commentaire de r\'esolution.
\end{steps}

\subsection{Escalader un signalement}

Si un signalement ne peut pas \^etre trait\'e localement, il peut \^etre escalad\'e\,:

\begin{steps}
  \item Ouvrir le signalement.
  \item Cliquer sur \textbf{Escalader vers les Services Centraux}.
  \item Le signalement est marqu\'e comme escalad\'e et remonte dans la liste des SC.
\end{steps}

%% ============================================================
\section{Journal d'audit}
%% ============================================================

\begin{rolebox}{services-centraux}
  Le journal d'audit est exclusivement accessible aux Services Centraux.
\end{rolebox}

Le journal d'audit enregistre toutes les actions effectu\'ees dans l'application (cr\'eation, modification, suppression d'utilisateurs, affectations, structures, etc.).

\screenshot{img_audit}{Journal d'audit}{Capturer la page du journal d'audit avec les filtres (type d'action, p\'eriode de dates, structure, utilisateur) et le tableau des entr\'ees d'audit (date/heure, action, utilisateur, donn\'ee concern\'ee, ancien/nouveau valeur).}

\subsection{Consulter le journal}

\begin{steps}
  \item Naviguer vers \textbf{Audit}.
  \item Appliquer les filtres souhaait\'es\,:
    \begin{itemize}
      \item \textit{Type d'action} : cr\'eation, modification, suppression.
      \item \textit{P\'eriode} : plage de dates.
      \item \textit{Structure} : limiter \`a une composante ou sous-structure.
      \item \textit{Utilisateur} : actions d'un utilisateur sp\'ecifique.
    \end{itemize}
  \item Parcourir les entr\'ees du journal.
\end{steps}

\subsection{Exporter le journal}

\begin{steps}
  \item Appliquer les filtres souhaait\'es pour r\'eduire le volume.
  \item Cliquer sur \textbf{Exporter (CSV)}.
  \item Le fichier CSV est t\'el\'echarg\'e avec les entr\'ees filtr\'ees.
\end{steps}

\chapter{Manuel --- Directeur de Composante et Directions Administratives}
\label{chap:dircomp}

\begin{rolebox}{directeur-composante, administrateur, directeur-administratif, directeur-administratif-adjoint}
  Ce chapitre d\'ecrit les fonctionnalit\'es des r\^oles de direction au niveau composante. Ces r\^oles disposent d'un large p\'erim\`etre de gestion, limit\'e cependant \`a leur \textbf{composante et ses sous-structures}.
\end{rolebox}

\section{Vue d'ensemble du profil}

Le directeur de composante (et les r\^oles assimil\'es) dispose des fonctionnalit\'es suivantes\,:

\begin{itemize}
  \item \textbf{Gestion des utilisateurs} dans son p\'erim\`etre (cr\'eation, modification, suppression).
  \item \textbf{Gestion des affectations} au sein de sa composante et de ses sous-structures.
  \item \textbf{Import et export} de donn\'ees pour sa composante.
  \item \textbf{G\'en\'eration d'organigrammes} pour sa composante et ses descendants.
  \item \textbf{Gestion des d\'el\'egations}.
  \item \textbf{Demande de r\^oles} personnalis\'es.
  \item \textbf{Signalements}.
\end{itemize}

\begin{notebox}[Diff\'erences entre les r\^oles]
  \begin{itemize}
    \item \texttt{directeur-composante} : acc\`es complet au niveau composante, y compris suppression d'utilisateurs et import/export.
    \item \texttt{administrateur} : droits administratifs g\'en\'eraux, similaires au directeur composante.
    \item \texttt{directeur-administratif} : gestion administrative, sans droit de suppression d'utilisateurs.
    \item \texttt{directeur-administratif-adjoint} : droits r\'eduits (pas de suppression, modification limit\'ee).
  \end{itemize}
\end{notebox}

\section{Tableau de bord}

\screenshot{img_dashboard_dircomp}{Tableau de bord -- Directeur de Composante}{Capturer le tableau de bord avec un compte directeur-composante. Montrer\,: l'ann\'ee courante (non modifiable), les tuiles de raccourcis disponibles pour ce profil, et la cloche de notifications.}

Le tableau de bord affiche\,:
\begin{itemize}
  \item L'\textbf{ann\'ee acad\'emique courante} (non modifiable par ce profil).
  \item Les \textbf{raccourcis} vers les fonctionnalit\'es accessibles.
  \item Les \textbf{notifications} en attente.
\end{itemize}

\section{Gestion des utilisateurs de la composante}

\begin{rolebox}{directeur-composante, administrateur, directeur-administratif}
  La gestion des utilisateurs est limit\'ee aux utilisateurs rattach\'es \`a la composante de l'utilisateur connect\'e.
\end{rolebox}

\subsection{Rechercher un utilisateur}

\begin{steps}
  \item Naviguer vers \textbf{Responsables}.
  \item Utiliser la barre de recherche (nom, pr\'enom, identifiant ou courriel).
  \item Les filtres hi\'erarchiques permettent de pr\'eciser\,: \textit{Composante}, \textit{D\'epartement}, \textit{Mention}, \textit{Parcours}, \textit{Niveau}.
  \item Les r\'esultats n'affichent que les utilisateurs du p\'erim\`etre autoris\'e.
\end{steps}

\screenshot{img_responsables_dircomp}{Page Responsables -- vue Directeur Composante}{Capturer la page Responsables avec les filtres hi\'erarchiques et la liste de r\'esultats. Montrer que la recherche est filtr\'ee au p\'erim\`etre de la composante.}

\subsection{Cr\'eer un utilisateur}

\begin{steps}
  \item Cliquer sur \textbf{Cr\'eer un utilisateur}.
  \item Remplir les champs obligatoires\,: \textit{Nom}, \textit{Pr\'enom}, \textit{Identifiant}.
  \item Remplir les champs optionnels\,: \textit{Courriel}, \textit{T\'el\'ephone}, \textit{Bureau}.
  \item Cliquer sur \textbf{Enregistrer}.
  \item Cr\'eer ensuite une affectation pour lier l'utilisateur \`a un r\^ole dans la composante (voir section~\ref{sec:dircomp_affectations}).
\end{steps}

\subsection{Modifier un utilisateur}

\begin{steps}
  \item Rechercher et ouvrir la fiche de l'utilisateur.
  \item Cliquer sur \textbf{Modifier}.
  \item Mettre \`a jour les informations autoris\'ees.
  \item Cliquer sur \textbf{Enregistrer}.
\end{steps}

\begin{warnbox}[Auto-modification interdite]
  Un directeur ne peut pas modifier sa propre fiche par ce biais. Pour modifier ses propres coordonn\'ees, utiliser la rubrique \textbf{Mon Profil}.
\end{warnbox}

\subsection{Supprimer un utilisateur}

\begin{rolebox}{directeur-composante, administrateur}
  La suppression d'utilisateurs est r\'eserv\'ee au directeur de composante et \`a l'administrateur.
\end{rolebox}

\begin{steps}
  \item Rechercher et ouvrir la fiche de l'utilisateur.
  \item Cliquer sur \textbf{Supprimer}.
  \item Confirmer la suppression.
\end{steps}

\begin{warnbox}
  La suppression d'un utilisateur entra\^ine la suppression de toutes ses affectations. Action irr\'eversible.
\end{warnbox}

\section{Gestion des affectations}
\label{sec:dircomp_affectations}

\subsection{Cr\'eer une affectation}

\begin{steps}
  \item Rechercher l'utilisateur concern\'e.
  \item Cliquer sur \textbf{Ajouter un r\^ole} dans sa fiche.
  \item S\'electionner le \textit{R\^ole}, la \textit{Composante}, et optionnellement les niveaux inf\'erieurs (\textit{D\'epartement}, \textit{Mention}, etc.).
  \item S\'electionner \'eventuellement un \textit{superviseur N+1}.
  \item Cliquer sur \textbf{Enregistrer}.
\end{steps}

\begin{notebox}
  Le directeur de composante ne peut cr\'eer des affectations que pour sa propre composante et ses sous-structures. Il ne peut pas affecter un utilisateur \`a une autre composante.
\end{notebox}

\screenshot{img_affectation_creer_dircomp}{Cr\'eation d'affectation -- Directeur Composante}{Capturer le formulaire de cr\'eation d'affectation avec les menus d\'eroulants (r\^ole, composante pr\'e-remplie, d\'epartement, mention, parcours, niveau, N+1).}

\subsection{G\'erer les contacts de r\^ole}

\begin{steps}
  \item Localiser l'affectation dans la liste.
  \item Cliquer sur \textbf{Contact de r\^ole}.
  \item Renseigner \textit{Courriel fonctionnel}, \textit{T\'el\'ephone de service}, \textit{Bureau fonctionnel}.
  \item Cliquer sur \textbf{Enregistrer}.
\end{steps}

\section{Import et Export}

\begin{rolebox}{directeur-composante, directeur-administratif, directeur-administratif-adjoint}
  L'import et l'export sont disponibles pour les r\^oles de direction de composante, limit\'es \`a leur p\'erim\`etre.
\end{rolebox}

\subsection{Exporter les donn\'ees de la composante}

\begin{steps}
  \item Naviguer vers \textbf{Import / Export}.
  \item Dans la section Export, s\'electionner\,:
    \begin{itemize}
      \item L'\textit{Ann\'ee} (toujours l'ann\'ee EN\_COURS pour ce profil).
      \item La \textit{Structure} (votre composante ou une sous-structure).
    \end{itemize}
  \item Cliquer sur \textbf{T\'el\'echarger}.
\end{steps}

\subsection{Importer des donn\'ees}

\begin{steps}
  \item Dans la section Import, d\'eposer le fichier \texttt{.xlsx}.
  \item S\'electionner la \textit{Structure cible} (votre composante).
  \item Cliquer sur \textbf{Pr\'evisualiser}.
  \item V\'erifier les op\'erations pr\'evues.
  \item Cliquer sur \textbf{Confirmer l'import}.
\end{steps}

\begin{warnbox}[P\'erim\`etre d'import]
  L'import ne peut modifier que les donn\'ees de la composante s\'electionn\'ee. Les donn\'ees d'autres composantes pr\'esentes dans le fichier sont ignor\'ees.
\end{warnbox}

\section{Organigrammes}

\subsection{G\'en\'erer un organigramme de la composante}

\begin{steps}
  \item Naviguer vers \textbf{Organigramme}.
  \item S\'electionner la vue (\textbf{Structures} ou \textbf{Personnes}).
  \item Choisir la \textit{structure racine} (votre composante ou une sous-structure).
  \item Appliquer les filtres souhait\'es.
  \item Cliquer sur \textbf{G\'en\'erer}.
\end{steps}

\screenshot{img_organigramme_dircomp}{G\'en\'eration d'organigramme -- Directeur Composante}{Capturer la page Organigramme avec les contr\^oles de g\'en\'eration\,: s\'electeur de vue (Structures/Personnes), s\'electeur de structure racine, filtres, et le bouton G\'en\'erer. Si un organigramme est d\'ej\`a affich\'e, montrer-le.}

\begin{notebox}[Consultation hors p\'erim\`etre]
  Un directeur de composante peut \textbf{consulter} des organigrammes d\'ej\`a g\'en\'er\'es pour d'autres structures (via la biblioth\`eque), mais il ne peut en g\'en\'erer que pour sa propre composante.
\end{notebox}

\section{D\'el\'egations}

\subsection{Cr\'eer une d\'el\'egation}

\begin{steps}
  \item Naviguer vers \textbf{D\'el\'egations}.
  \item Cliquer sur \textbf{Cr\'eer une d\'el\'egation}.
  \item S\'electionner le \textit{D\'el\'egu\'e}, la \textit{Structure concern\'ee} et les \textit{dates}.
  \item Cliquer sur \textbf{Enregistrer}.
\end{steps}

\screenshot{img_delegations_dircomp}{Gestion des d\'el\'egations -- Directeur Composante}{Capturer la page D\'el\'egations avec la liste des d\'el\'egations en cours et le bouton de cr\'eation.}

\section{Demandes de r\^oles}

Le directeur de composante peut soumettre des demandes de r\^oles personnalis\'es pour des membres de sa composante.

\subsection{Soumettre une demande}

\begin{steps}
  \item Naviguer vers \textbf{Demandes de r\^oles}.
  \item Cliquer sur \textbf{Nouvelle demande}.
  \item Remplir\,: \textit{Utilisateur concern\'e}, \textit{R\^ole demand\'e}, \textit{Structure}, \textit{Justification}.
  \item Cliquer sur \textbf{Envoyer}.
  \item La demande est transmise aux Services Centraux pour examen.
\end{steps}

\screenshot{img_demandes_roles_dircomp}{Demandes de r\^oles -- Directeur Composante}{Capturer la page demandes de r\^oles avec le formulaire de soumission et la liste des demandes existantes (en attente, approuv\'ees, refus\'ees).}

\subsection{Suivre l'\'etat d'une demande}

\begin{itemize}
  \item \textbf{En attente} : la demande est chez les Services Centraux.
  \item \textbf{Approuv\'ee} : le r\^ole a \'et\'e attribu\'e \`a l'utilisateur.
  \item \textbf{Refus\'ee} : la demande a \'et\'e rejet\'ee (un commentaire de justification peut \^etre pr\'esent).
\end{itemize}

\begin{notebox}
  Une notification est envoy\'ee automatiquement lorsque la demande est trait\'ee par les Services Centraux.
\end{notebox}

\section{Signalements}

\begin{steps}
  \item Naviguer vers \textbf{Signalements}.
  \item Cliquer sur \textbf{Cr\'eer un signalement}.
  \item S\'electionner le \textit{type d'erreur} (nom, pr\'enom, t\'el\'ephone, bureau, courriel, r\^ole, autre).
  \item Identifier la \textit{personne} ou la \textit{structure} concern\'ee.
  \item D\'ecrire le probl\`eme dans le champ \textit{Description}.
  \item Cliquer sur \textbf{Envoyer}.
\end{steps}

\screenshot{img_signalements_dircomp}{Signalements -- Directeur Composante}{Capturer la page Signalements avec le formulaire de cr\'eation et la liste des signalements existants. Montrer les filtres de statut.}

\chapter{Manuel --- Directions et Responsables P\'edagogiques}
\label{chap:directions}

\begin{rolebox}{directeur-departement, vice-president-departement, directeur-adjoint-licence, responsable-service-pedagogique, responsable-adjoint-service-pedagogique, directeur-mention, directeur-specialite, responsable-formation, responsable-annee, directeur-etudes, responsable-qualite, responsable-international, referent-commun, directeur-adjoint-ecole}
  Ce chapitre couvre l'ensemble des r\^oles de direction au niveau d\'epartement, mention, parcours, niveau ou service. Ces r\^oles ont des droits d'\'ecriture limit\'es \`a leur p\'erim\`etre et ne peuvent pas g\'erer les utilisateurs.
\end{rolebox}

\section{Vue d'ensemble du profil}

Les r\^oles de direction et de responsabilit\'e p\'edagogique disposent des fonctionnalit\'es suivantes\,:

\begin{itemize}
  \item \textbf{Consultation de l'annuaire} (responsables, formations, structures, secr\'etariats).
  \item \textbf{Gestion des affectations} dans leur p\'erim\`etre (cr\'eation et modification de contacts de r\^ole).
  \item \textbf{G\'en\'eration d'organigrammes} pour leur structure et ses descendants.
  \item \textbf{Consultation de la biblioth\`eque} d'organigrammes.
  \item \textbf{Gestion des d\'el\'egations}.
  \item \textbf{Demandes de r\^oles} personnalis\'es.
  \item \textbf{Signalements}.
\end{itemize}

\begin{notebox}[Ce qui n'est PAS disponible]
  Ces r\^oles \textbf{ne peuvent pas}\,:
  \begin{itemize}
    \item Cr\'eer, modifier ou supprimer des comptes utilisateurs.
    \item Importer ou exporter des donn\'ees (sauf si la d\'el\'egation le pr\'evoit).
    \item G\'erer les ann\'ees acad\'emiques.
    \item Modifier les structures.
    \item Consulter le journal d'audit.
  \end{itemize}
\end{notebox}

\section{Tableau de bord}

\screenshot{img_dashboard_direction}{Tableau de bord -- Directions et Responsables}{Capturer le tableau de bord avec un compte de type direction (ex. directeur-departement ou directeur-mention). Montrer les tuiles disponibles pour ce profil, l'ann\'ee courante (non modifiable) et les notifications.}

\section{Annuaire -- Recherche}

L'annuaire est la fonctionnalit\'e principale pour les r\^oles de direction. Il permet de retrouver rapidement les responsables, formations, structures et secr\'etariats.

\begin{steps}
  \item Naviguer vers \textbf{Rechercher} dans le menu.
  \item S\'electionner l'onglet souhait\'e\,: \textbf{Responsables}, \textbf{Formations}, \textbf{Structures} ou \textbf{Secr\'etariats}.
  \item Saisir un texte dans la barre de recherche et/ou utiliser les filtres.
  \item Les r\'esultats s'affichent en dessous.
\end{steps}

\screenshot{img_annuaire_direction}{Annuaire -- Vue Directions}{Capturer la page Recherche avec l'onglet Responsables actif, une recherche effectu\'ee et les r\'esultats visibles. Montrer les filtres (composante, d\'epartement, mention) et le tableau de r\'esultats.}

Pour les d\'etails complets de l'annuaire, voir le chapitre~\ref{chap:commun}.

\section{Gestion des affectations dans le p\'erim\`etre}

Les r\^oles de direction peuvent g\'erer les affectations au sein de leur structure.

\subsection{Cr\'eer ou modifier une affectation}

\begin{steps}
  \item Naviguer vers \textbf{Responsables}.
  \item Rechercher la personne concern\'ee.
  \item Cliquer sur \textbf{Ajouter un r\^ole} (dans la fiche de la personne).
  \item S\'electionner le \textit{R\^ole}, la \textit{Structure} (limit\'ee au p\'erim\`etre autoris\'e).
  \item Cliquer sur \textbf{Enregistrer}.
\end{steps}

\begin{notebox}[P\'erim\`etre des affectations]
  Un directeur de d\'epartement ne peut cr\'eer des affectations que pour son d\'epartement et ses sous-structures (mentions, parcours, niveaux). Il ne peut pas affecter des r\^oles dans d'autres d\'epartements.
\end{notebox}

\screenshot{img_affectation_direction}{Cr\'eation d'affectation -- Direction}{Capturer le formulaire d'affectation pour un r\^ole de direction. Montrer que les menus d\'eroulants de structure sont limit\'es au p\'erim\`etre du r\^ole (ex. d\'epartements disponibles = uniquement le d\'epartement de l'utilisateur).}

\subsection{Mettre \`a jour un contact de r\^ole}

\begin{steps}
  \item Localiser l'affectation dans la liste \textbf{Responsables}.
  \item Cliquer sur \textbf{Contact de r\^ole}.
  \item Renseigner ou mettre \`a jour\,: \textit{Courriel fonctionnel}, \textit{T\'el\'ephone}, \textit{Bureau}.
  \item Cliquer sur \textbf{Enregistrer}.
\end{steps}

\begin{tipbox}
  La mise \`a jour des contacts de r\^ole est particuli\`erement importante pour les secr\'etariats et les responsables de formation, car ces informations sont visibles dans l'annuaire par tous les utilisateurs.
\end{tipbox}

\section{Fiches Structures}

Les directions peuvent consulter les fiches de structure dans leur p\'erim\`etre.

\begin{steps}
  \item Naviguer vers \textbf{Structures} ou rechercher dans l'annuaire.
  \item Cliquer sur la structure souhait\'ee.
  \item La fiche d\'etaill\'ee affiche\,: hi\'erarchie, responsables, secr\'etariat, sous-structures, statistiques.
\end{steps}

\screenshot{img_fiche_structure_direction}{Fiche structure -- vue Direction}{Capturer la fiche d\'etaill\'ee d'une structure (d\'epartement ou mention). Montrer les sections\,: titre, hi\'erarchie parente, liste des responsables avec leurs r\^oles, liste du secr\'etariat, sous-structures, et statistiques chiffr\'ees.}

\section{Organigrammes}

\subsection{G\'en\'erer un organigramme pour son p\'erim\`etre}

\begin{steps}
  \item Naviguer vers \textbf{Organigramme}.
  \item Choisir la vue (\textbf{Structures} ou \textbf{Personnes}).
  \item S\'electionner la structure racine parmi celles de son p\'erim\`etre.
  \item Appliquer les filtres (r\^ole, niveau, etc.).
  \item Cliquer sur \textbf{G\'en\'erer}.
\end{steps}

\screenshot{img_organigramme_direction}{Organigramme -- G\'en\'eration pour une Direction}{Capturer la page organigramme avec la vue Personnes, le s\'electeur de structure racine (montrant le d\'epartement ou la mention de l'utilisateur) et les filtres. Montrer le r\'esultat g\'en\'er\'e si possible.}

\subsection{Consulter la biblioth\`eque}

\begin{steps}
  \item Dans la page \textbf{Organigramme}, acc\'eder \`a la section \og Biblioth\`eque \fg{}.
  \item Filtrer par structure ou type de vue.
  \item Cliquer sur \textbf{Voir en structures} ou \textbf{Voir en personnes} pour consulter un organigramme.
\end{steps}

\begin{notebox}[Consultation \'etendue]
  Les r\^oles de direction peuvent consulter les organigrammes d\'ej\`a g\'en\'er\'es par d'autres utilisateurs, m\^eme pour des structures hors de leur p\'erim\`etre. En revanche, ils ne peuvent g\'en\'erer des organigrammes que pour leur p\'erim\`etre.
\end{notebox}

\section{D\'el\'egations}

Les directions peuvent d\'el\'eguer leurs droits \`a un autre utilisateur dans leur p\'erim\`etre.

\subsection{Cr\'eer une d\'el\'egation}

\begin{steps}
  \item Naviguer vers \textbf{D\'el\'egations}.
  \item Cliquer sur \textbf{Cr\'eer une d\'el\'egation}.
  \item S\'electionner le \textit{D\'el\'egu\'e} (utilisateur existant).
  \item S\'electionner la \textit{Structure concern\'ee} (dans votre p\'erim\`etre).
  \item Indiquer les \textit{dates de d\'ebut et de fin}.
  \item Cliquer sur \textbf{Enregistrer}.
\end{steps}

\screenshot{img_delegations_direction}{D\'el\'egations -- Directions}{Capturer la page D\'el\'egations avec la liste des d\'el\'egations actives et le formulaire de cr\'eation visible (ou le bouton Cr\'eer).}

\subsection{R\'evoquer une d\'el\'egation}

\begin{steps}
  \item Dans la liste, cliquer sur la d\'el\'egation.
  \item Cliquer sur \textbf{R\'evoquer}.
  \item Confirmer.
\end{steps}

\section{Demandes de r\^oles}

Les directions peuvent soumettre des demandes de r\^oles personnalis\'es pour des personnes de leur p\'erim\`etre.

\subsection{Soumettre une demande}

\begin{steps}
  \item Naviguer vers \textbf{Demandes de r\^oles}.
  \item Cliquer sur \textbf{Nouvelle demande}.
  \item Renseigner\,: \textit{Utilisateur}, \textit{R\^ole demand\'e}, \textit{Structure}, \textit{Justification}.
  \item Cliquer sur \textbf{Envoyer}.
\end{steps}

\screenshot{img_demandes_roles_direction}{Demandes de r\^oles -- Directions}{Capturer la page demandes de r\^oles avec la liste des demandes (en attente, approuv\'ees, refus\'ees) et le formulaire ou bouton de nouvelle demande.}

\subsection{Suivre ses demandes}

Les statuts possibles sont\,:
\begin{itemize}
  \item \textbf{En attente} : examen par les Services Centraux en cours.
  \item \textbf{Approuv\'ee} : r\^ole attribu\'e, affectation cr\'e\'ee.
  \item \textbf{Refus\'ee} : demande rejet\'ee (consulter le commentaire de justification).
\end{itemize}

\section{Signalements}

\begin{steps}
  \item Naviguer vers \textbf{Signalements}.
  \item Cliquer sur \textbf{Cr\'eer un signalement}.
  \item S\'electionner le \textit{type d'erreur} et la \textit{personne ou structure concern\'ee}.
  \item D\'ecrire le probl\`eme.
  \item Cliquer sur \textbf{Envoyer}.
\end{steps}

\screenshot{img_signalements_direction}{Signalements -- Directions}{Capturer la page Signalements. Montrer le formulaire de cr\'eation et la liste des signalements de l'utilisateur avec les statuts.}

\chapter{Manuel --- Secr\'etariat P\'edagogique}
\label{chap:secretariat}

\begin{rolebox}{secretariat-pedagogique}
  Ce chapitre est destin\'e aux agents du secr\'etariat p\'edagogique. Ce profil est orient\'e \textbf{consultation et mise \`a jour des contacts fonctionnels}. Il ne permet pas de g\'erer des utilisateurs, de cr\'eer des affectations ou de g\'en\'erer des organigrammes.
\end{rolebox}

\section{Vue d'ensemble du profil}

Le profil \texttt{secretariat-pedagogique} dispose des fonctionnalit\'es suivantes\,:

\begin{itemize}
  \item \textbf{Consultation de l'annuaire} : recherche de responsables, formations, structures et secr\'etariats.
  \item \textbf{Consultation des fiches structures} : acc\`es aux d\'etails complets des structures.
  \item \textbf{Mise \`a jour des contacts de r\^ole} : modification des informations de contact fonctionnel li\'ees aux affectations (courriel fonctionnel, t\'el\'ephone, bureau).
  \item \textbf{Consultation des organigrammes} g\'en\'er\'es (lecture seule).
  \item \textbf{Signalements} : signalement d'erreurs sur les donn\'ees.
  \item \textbf{Modification du profil personnel}.
\end{itemize}

\begin{notebox}[Ce qui n'est PAS disponible]
  Le secr\'etariat p\'edagogique ne peut \textbf{pas}\,:
  \begin{itemize}
    \item Cr\'eer, modifier ou supprimer des comptes utilisateurs.
    \item Cr\'eer ou supprimer des affectations.
    \item G\'en\'erer de nouveaux organigrammes.
    \item Importer ou exporter des donn\'ees.
    \item G\'erer des d\'el\'egations.
    \item G\'erer les ann\'ees acad\'emiques.
  \end{itemize}
\end{notebox}

\section{Tableau de bord}

\screenshot{img_dashboard_secretariat}{Tableau de bord -- Secr\'etariat P\'edagogique}{Capturer le tableau de bord apr\`es connexion avec un compte secretariat-pedagogique. Montrer les tuiles de raccourcis limit\'ees \`a ce profil (Rechercher, Organigramme, Signalements, Profil) et l'ann\'ee courante.}

\section{Consultation de l'annuaire}

L'annuaire est la fonctionnalit\'e principale pour le secr\'etariat. Il permet de retrouver rapidement toute personne, formation ou structure de l'\'etablissement.

\begin{steps}
  \item Cliquer sur \textbf{Rechercher} dans le menu.
  \item Choisir l'onglet correspondant \`a votre recherche.
  \item Saisir un texte libre ou utiliser les filtres.
\end{steps}

\subsection{Onglet Responsables}

Permet de trouver une personne par\,: nom, pr\'enom, courriel, r\^ole, composante, d\'epartement, mention, parcours ou niveau.

\screenshot{img_annuaire_responsables_secretariat}{Annuaire -- Onglet Responsables (Secr\'etariat)}{Capturer la page Recherche avec l'onglet Responsables actif. Montrer la barre de recherche texte, les filtres hi\'erarchiques (composante, d\'epartement, mention, parcours, niveau) et les filtres de r\^ole, et un exemple de r\'esultats avec des cartes ou lignes de r\'esultat.}

\subsection{Onglet Formations}

Permet de rechercher une formation par son intitul\'e, son type de dipl\^ome, sa composante ou son d\'epartement.

\screenshot{img_annuaire_formations_secretariat}{Annuaire -- Onglet Formations (Secr\'etariat)}{Capturer l'onglet Formations de la recherche avec des r\'esultats affich\'es. Montrer les filtres disponibles (type de dipl\^ome, composante, d\'epartement) et les cartes de r\'esultat avec intitul\'e de formation et responsables.}

\subsection{Onglet Structures}

Permet de retrouver une structure (composante, d\'epartement, mention, parcours, service) par son nom ou sa hi\'erarchie.

\screenshot{img_annuaire_structures_secretariat}{Annuaire -- Onglet Structures (Secr\'etariat)}{Capturer l'onglet Structures de l'annuaire avec des r\'esultats. Montrer les filtres et le tableau de r\'esultats (nom de structure, type, parent).}

\subsection{Onglet Secr\'etariats}

Permet de retrouver les contacts de secr\'etariat d'une composante ou d'un d\'epartement.

\screenshot{img_annuaire_secretariats_secretariat}{Annuaire -- Onglet Secr\'etariats (Secr\'etariat)}{Capturer l'onglet Secr\'etariats avec des r\'esultats. Montrer les informations affich\'ees\,: nom de la personne, structure rattach\'ee, courriel fonctionnel, t\'el\'ephone, bureau.}

\section{Consultation des fiches structures}

Le secr\'etariat peut consulter les fiches d\'etaill\'ees de toutes les structures.

\begin{steps}
  \item Rechercher la structure dans l'annuaire (onglet \textbf{Structures}).
  \item Cliquer sur la structure pour ouvrir sa fiche.
  \item La fiche affiche\,: identification, hi\'erarchie, coordonn\'ees, responsables, secr\'etariat, sous-structures.
\end{steps}

\screenshot{img_fiche_structure_secretariat}{Fiche structure -- Vue Secr\'etariat}{Capturer la fiche d\'etaill\'ee d'une composante ou d'un d\'epartement. Toutes les sections doivent \^etre visibles\,: titre, type, ann\'ee, parent, coordonn\'ees, liste des responsables (nom, r\^ole, courriel), liste du secr\'etariat, sous-structures directes.}

\section{Mise \`a jour des contacts de r\^ole}

Le secr\'etariat p\'edagogique peut mettre \`a jour les informations de contact fonctionnel associ\'ees aux affectations (courriel fonctionnel, t\'el\'ephone de service, bureau).

\begin{warnbox}[Distinction contact personnel / contact fonctionnel]
  Le \textbf{contact personnel} (courriel priv\'e, t\'el\'ephone personnel) est g\'er\'e via \og Mon Profil \fg{}. \\
  Le \textbf{contact de r\^ole} (courriel fonctionnel de service, t\'el\'ephone du secr\'etariat, bureau du service) est g\'er\'e via la rubrique \textbf{Responsables} \`a travers les affectations.
\end{warnbox}

\subsection{Trouver l'affectation \`a mettre \`a jour}

\begin{steps}
  \item Naviguer vers \textbf{Responsables}.
  \item Rechercher la personne ou filtrer par structure.
  \item Localiser l'affectation dont les coordonn\'ees doivent \^etre mises \`a jour.
\end{steps}

\subsection{Modifier le contact de r\^ole}

\begin{steps}
  \item Dans la ligne ou fiche de l'affectation, cliquer sur \textbf{Contact de r\^ole} (ou ic\^one crayon).
  \item Modifier les champs souhait\'es\,:
    \begin{itemize}
      \item \textit{Courriel fonctionnel} : courriel de service (ex. : \texttt{secretariat.droit@univ.fr}).
      \item \textit{T\'el\'ephone de service} : num\'ero du bureau ou du secr\'etariat.
      \item \textit{Bureau / Localisation} : num\'ero de bureau, b\^atiment, campus.
    \end{itemize}
  \item Cliquer sur \textbf{Enregistrer}.
\end{steps}

\screenshot{img_contact_role_secretariat}{Modification du contact de r\^ole -- Secr\'etariat}{Capturer le formulaire d'\'edition du contact de r\^ole tel qu'il appara\^it pour un secr\'etariat. Montrer les trois champs (courriel fonctionnel, t\'el\'ephone, bureau) avec leurs labels et un exemple de valeur saisie.}

\begin{tipbox}[Cas d'usage typique]
  En d\'ebut d'ann\'ee acad\'emique ou lors d'un changement de locaux, le secr\'etariat met \`a jour les coordonn\'ees fonctionnelles de toutes les affectations de son service afin que l'annuaire reste \`a jour pour l'ensemble des utilisateurs.
\end{tipbox}

\section{Consultation des organigrammes}

Le secr\'etariat peut consulter les organigrammes d\'ej\`a g\'en\'er\'es mais ne peut pas en cr\'eer de nouveaux.

\begin{steps}
  \item Naviguer vers \textbf{Organigramme}.
  \item Acc\'eder \`a la section \og Biblioth\`eque \fg{}.
  \item Filtrer par structure ou type de vue.
  \item Cliquer sur \textbf{Voir en structures} ou \textbf{Voir en personnes}.
\end{steps}

\screenshot{img_organigramme_consultation_secretariat}{Consultation d'organigramme -- Secr\'etariat}{Capturer l'organigramme en consultation (vue personnes ou vue structures). Montrer que les boutons de g\'en\'eration et de gel ne sont pas visibles (ou sont d\'esactiv\'es) pour ce profil.}

\begin{notebox}
  Le secr\'etariat peut consulter les organigrammes de toutes les structures, pas uniquement ceux de sa propre composante.
\end{notebox}

\section{Signalements}

Le secr\'etariat peut signaler des erreurs sur les donn\'ees.

\begin{steps}
  \item Naviguer vers \textbf{Signalements}.
  \item Cliquer sur \textbf{Cr\'eer un signalement}.
  \item S\'electionner le type d'erreur\,:
    \begin{itemize}
      \item \textbf{Nom} : erreur sur le nom d'une personne.
      \item \textbf{Pr\'enom} : erreur sur le pr\'enom.
      \item \textbf{T\'el\'ephone} : num\'ero incorrect ou manquant.
      \item \textbf{Bureau} : localisation incorrecte.
      \item \textbf{Courriel} : adresse courriel incorrecte.
      \item \textbf{R\^ole} : affectation incorrecte ou manquante.
      \item \textbf{Autre} : tout autre type d'erreur.
    \end{itemize}
  \item Identifier la \textit{personne} ou la \textit{structure} concern\'ee.
  \item D\'ecrire le probl\`eme dans le champ \textit{Description}.
  \item Cliquer sur \textbf{Envoyer}.
\end{steps}

\screenshot{img_signalements_secretariat}{Signalements -- Secr\'etariat P\'edagogique}{Capturer la page Signalements avec\,: (1) le formulaire de cr\'eation (menu d\'eroulant du type d'erreur, champ de recherche de la personne concern\'ee, champ description), et (2) la liste des signalements d\'ej\`a soumis avec leurs statuts.}

\begin{tipbox}[Suivi des signalements]
  Apr\`es soumission, vous pouvez suivre le statut de vos signalements depuis la page \textbf{Signalements}. Une notification vous est envoy\'ee lorsque votre signalement est trait\'e ou cl\^otur\'e.
\end{tipbox}

\section{Mon Profil}

Le secr\'etariat peut consulter et mettre \`a jour ses propres coordonn\'ees.

\begin{steps}
  \item Cliquer sur \textbf{Mon Profil} dans le menu.
  \item Cliquer sur \textbf{Modifier} dans la section coordonn\'ees.
  \item Mettre \`a jour\,: \textit{Courriel secondaire}, \textit{T\'el\'ephone}, \textit{Bureau}.
  \item Cliquer sur \textbf{Enregistrer}.
\end{steps}

\chapter{Manuel --- Utilisateur Standard et Lecture Seule}
\label{chap:utilisateur}

\begin{rolebox}{utilisateur-simple, lecture-seule}
  Ce chapitre s'adresse aux utilisateurs avec un acc\`es minimal. Ces profils sont orient\'es \textbf{consultation} et ne permettent aucune modification des donn\'ees (sauf les signalements et la mise \`a jour du profil personnel pour \texttt{utilisateur-simple}).
\end{rolebox}

\section{Diff\'erence entre les deux profils}

\medskip
\begin{tabularx}{\textwidth}{L{5cm}C{3cm}C{3cm}}
  \toprule
  \rowcolor{cgplightblue}
  \textbf{Fonctionnalit\'e} & \textbf{utilisateur-simple} & \textbf{lecture-seule} \\
  \midrule
  Consulter l'annuaire                     & {\color{cgpgreen}\checkmark} & {\color{cgpgreen}\checkmark} \\
  \rowcolor{cgplightgray}
  Consulter les organigrammes              & {\color{cgpgreen}\checkmark} & {\color{cgpgreen}\checkmark} \\
  Consulter les fiches structures          & {\color{cgpgreen}\checkmark} & {\color{cgpgreen}\checkmark} \\
  \rowcolor{cgplightgray}
  Cr\'eer des signalements                 & {\color{cgpgreen}\checkmark} & {\color{cgpgreen}\checkmark} \\
  Modifier son profil (coordonn\'ees)      & {\color{cgpgreen}\checkmark} & {\color{cgpred}$\times$} \\
  \rowcolor{cgplightgray}
  G\'en\'erer des organigrammes            & {\color{cgpred}$\times$}    & {\color{cgpred}$\times$} \\
  G\'erer des affectations / utilisateurs  & {\color{cgpred}$\times$}    & {\color{cgpred}$\times$} \\
  \bottomrule
\end{tabularx}

\section{Tableau de bord}

\screenshot{img_dashboard_utilisateur}{Tableau de bord -- Utilisateur Standard}{Capturer le tableau de bord avec un compte utilisateur-simple ou lecture-seule. Montrer les tuiles de raccourcis limit\'ees \`a ce profil (Rechercher, Organigramme, Signalements, Profil). L'absence de tuiles de gestion (Responsables, Ann\'ees, Import/Export) doit \^etre apparente.}

\section{Consultation de l'annuaire}

L'annuaire est la fonctionnalit\'e principale des utilisateurs standard. Il permet de retrouver rapidement n'importe quelle personne, formation ou structure de l'\'etablissement.

\subsection{Acc\`eder \`a l'annuaire}

\begin{steps}
  \item Cliquer sur \textbf{Rechercher} dans le menu de navigation.
  \item La page s'ouvre sur l'onglet \textbf{Responsables} par d\'efaut.
  \item Quatre onglets sont disponibles\,: \textbf{Responsables}, \textbf{Formations}, \textbf{Structures}, \textbf{Secr\'etariats}.
\end{steps}

\screenshot{img_annuaire_utilisateur}{Annuaire -- Vue Utilisateur Standard}{Capturer la page de recherche avec les quatre onglets visibles en haut. Montrer l'onglet Responsables actif avec une recherche effectu\'ee et des r\'esultats affich\'es. La barre de recherche et les filtres principaux doivent \^etre visibles.}

\subsection{Rechercher un responsable}

\begin{steps}
  \item S\'electionner l'onglet \textbf{Responsables}.
  \item Saisir un nom, pr\'enom ou intitul\'e dans la barre de recherche.
  \item Affiner avec les filtres\,: \textit{Composante}, \textit{D\'epartement}, \textit{Mention}, \textit{Parcours}, \textit{Niveau}, \textit{R\^ole}.
  \item Les r\'esultats s'affichent sous forme de cartes ou de tableau avec\,: nom, pr\'enom, r\^ole, structure, courriel institutionnel.
\end{steps}

\begin{tipbox}[Conseils de recherche]
  \begin{itemize}
    \item Pour une recherche rapide, saisissez le nom de famille dans la barre de recherche.
    \item Pour trouver tous les responsables d'un d\'epartement, utilisez uniquement le filtre \textit{D\'epartement} sans texte de recherche.
    \item Pour trouver tous les directeurs de mention, utilisez le filtre \textit{R\^ole} = \og Directeur de mention \fg{}.
  \end{itemize}
\end{tipbox}

\subsection{Rechercher une formation}

\begin{steps}
  \item S\'electionner l'onglet \textbf{Formations}.
  \item Saisir le nom de la formation ou utiliser les filtres\,: \textit{Type de dipl\^ome}, \textit{Composante}, \textit{D\'epartement}.
  \item Les r\'esultats affichent\,: intitul\'e de la formation, type de dipl\^ome, composante, directeur de mention, directeur d'\'etudes.
\end{steps}

\screenshot{img_annuaire_formations_utilisateur}{Annuaire Formations -- Utilisateur Standard}{Capturer l'onglet Formations avec des r\'esultats. Montrer les filtres disponibles et les cartes de r\'esultats avec l'intitul\'e de la formation, le type de dipl\^ome (Licence, Master, etc.) et les responsables associ\'es.}

\subsection{Rechercher une structure}

\begin{steps}
  \item S\'electionner l'onglet \textbf{Structures}.
  \item Saisir le nom de la structure ou utiliser les filtres hi\'erarchiques.
  \item Cliquer sur une structure dans les r\'esultats pour afficher sa fiche d\'etaill\'ee.
\end{steps}

\subsection{Rechercher un secr\'etariat}

\begin{steps}
  \item S\'electionner l'onglet \textbf{Secr\'etariats}.
  \item Filtrer par \textit{Composante} ou \textit{D\'epartement}.
  \item Les r\'esultats affichent les contacts du secr\'etariat\,: nom, courriel fonctionnel, t\'el\'ephone, bureau.
\end{steps}

\section{Consultation des fiches structures}

\begin{steps}
  \item Rechercher la structure souhait\'ee dans l'onglet \textbf{Structures} de l'annuaire.
  \item Cliquer sur la structure pour ouvrir sa fiche.
  \item La fiche affiche\,: identification, hi\'erarchie, coordonn\'ees, liste des responsables, secr\'etariat, sous-structures.
\end{steps}

\screenshot{img_fiche_structure_utilisateur}{Fiche structure -- Vue Utilisateur}{Capturer une fiche structure compl\`ete (composante ou d\'epartement). Toutes les sections doivent \^etre visibles. Montrer que les boutons de modification ne sont pas accessibles pour cet utilisateur.}

\section{Consultation des organigrammes}

L'utilisateur standard peut consulter les organigrammes existants mais ne peut pas en g\'en\'erer de nouveaux.

\begin{steps}
  \item Naviguer vers \textbf{Organigramme}.
  \item Dans la section \og Biblioth\`eque \fg{}, utiliser les filtres pour retrouver un organigramme.
  \item Cliquer sur \textbf{Voir en structures} ou \textbf{Voir en personnes}.
\end{steps}

\screenshot{img_organigramme_utilisateur}{Organigramme -- Consultation Utilisateur}{Capturer la page organigramme telle que vue par un utilisateur standard. Montrer la biblioth\`eque avec les organigrammes disponibles, et un organigramme ouvert en vue structures ou personnes. Les boutons de g\'en\'eration et de gel doivent \^etre absents.}

\begin{notebox}
  Si aucun organigramme n'a encore \'et\'e g\'en\'er\'e par un utilisateur autoris\'e, la biblioth\`eque sera vide. Dans ce cas, contactez votre directeur de composante ou les Services Centraux.
\end{notebox}

\section{Signalements}

\begin{rolebox}{utilisateur-simple, lecture-seule}
  Les deux profils peuvent soumettre des signalements.
\end{rolebox}

\begin{steps}
  \item Naviguer vers \textbf{Signalements}.
  \item Cliquer sur \textbf{Cr\'eer un signalement}.
  \item S\'electionner le type d'erreur (nom, pr\'enom, t\'el\'ephone, bureau, courriel, r\^ole, autre).
  \item Identifier la personne ou la structure concern\'ee.
  \item D\'ecrire clairement l'erreur constat\'ee.
  \item Cliquer sur \textbf{Envoyer}.
\end{steps}

\screenshot{img_signalements_utilisateur}{Signalements -- Utilisateur Standard}{Capturer la page Signalements avec le formulaire de cr\'eation (type d'erreur, personne concern\'ee, description) et la liste des signalements pr\'ec\'edemment soumis. Montrer les badges de statut (En attente, Trait\'e, Cl\^otur\'e).}

\begin{tipbox}[Bon usage des signalements]
  \begin{itemize}
    \item Soyez pr\'ecis dans la description\,: indiquez l'erreur constat\'ee et la valeur correcte si vous la connaissez.
    \item Ne soumettez pas plusieurs signalements pour la m\^eme erreur.
    \item Vous serez notifi\'e lorsque votre signalement sera trait\'e.
  \end{itemize}
\end{tipbox}

\section{Mon Profil}

\begin{rolebox}{utilisateur-simple}
  L'utilisateur simple peut modifier ses coordonn\'ees. L'utilisateur lecture-seule peut uniquement consulter son profil.
\end{rolebox}

\begin{steps}
  \item Cliquer sur \textbf{Mon Profil} dans le menu.
  \item Consulter les informations affich\'ees\,: nom, pr\'enom, identifiant, r\^oles, courriel, t\'el\'ephone, bureau.
  \item Pour \texttt{utilisateur-simple}\,: cliquer sur \textbf{Modifier} pour mettre \`a jour les coordonn\'ees.
  \item Modifier \textit{Courriel secondaire}, \textit{T\'el\'ephone} et/ou \textit{Bureau}.
  \item Cliquer sur \textbf{Enregistrer}.
\end{steps}

\screenshot{img_profil_utilisateur}{Mon Profil -- Utilisateur Standard}{Capturer la page Mon Profil d'un utilisateur simple. Montrer toutes les informations affich\'ees (identit\'e, r\^ole(s), coordonn\'ees). Le bouton Modifier doit \^etre visible. Utiliser des donn\'ees de test non confidentielles.}

\chapter{Fonctionnalit\'es Communes \`a Tous les Profils}
\label{chap:commun}

\begin{rolebox}{Tous les r\^oles}
  Ce chapitre d\'ecrit les fonctionnalit\'es accessibles \`a tous les utilisateurs, quel que soit leur r\^ole. Il compl\`ete les chapitres d\'edi\'es \`a chaque profil.
\end{rolebox}

\section{Annuaire -- Guide d\'etaill\'e}
\label{sec:annuaire}

L'annuaire est la fonctionnalit\'e de consultation centrale de CGP. Il est accessible depuis la rubrique \textbf{Rechercher} du menu.

\subsection{Pr\'esentation g\'en\'erale}

L'annuaire est organis\'e en quatre onglets\,:

\medskip
\begin{tabularx}{\textwidth}{L{3cm}X}
  \toprule
  \rowcolor{cgplightblue}
  \textbf{Onglet} & \textbf{Contenu} \\
  \midrule
  \textbf{Responsables} & Liste des personnes affect\'ees \`a un r\^ole dans une structure. \\
  \rowcolor{cgplightgray}
  \textbf{Formations} & Liste des formations (Licence, Master, etc.) avec leurs responsables. \\
  \textbf{Structures} & Liste des structures de l'\'etablissement (composantes, d\'epartements, mentions...). \\
  \rowcolor{cgplightgray}
  \textbf{Secr\'etariats} & Contacts fonctionnels du secr\'etariat p\'edagogique. \\
  \bottomrule
\end{tabularx}

\screenshot{img_annuaire_onglets}{Annuaire -- Vue g\'en\'erale avec les quatre onglets}{Capturer la page Recherche avec les quatre onglets visibles en haut (Responsables, Formations, Structures, Secr\'etariats). Montrer l'onglet Responsables actif avec une recherche et des r\'esultats. La barre de recherche et les filtres doivent \^etre visibles.}

\subsection{Onglet Responsables}

Cet onglet permet de rechercher toute personne ayant une affectation active dans l'ann\'ee courante.

\subsubsection{Filtres disponibles}

\begin{itemize}
  \item \textbf{Recherche libre} : saisie de texte (nom, pr\'enom, courriel, identifiant).
  \item \textbf{Composante} : filtre par composante (niveau 1 de la hi\'erarchie).
  \item \textbf{D\'epartement} : se d\'eroule apr\`es s\'election d'une composante.
  \item \textbf{Mention} : se d\'eroule apr\`es s\'election d'un d\'epartement.
  \item \textbf{Parcours} : se d\'eroule apr\`es s\'election d'une mention.
  \item \textbf{Niveau} : se d\'eroule apr\`es s\'election d'un parcours.
  \item \textbf{R\^ole} : filtre par type de r\^ole.
\end{itemize}

\begin{notebox}[Filtres en cascade]
  Les filtres hi\'erarchiques fonctionnent en cascade\,: la s\'election d'une composante filtre automatiquement les d\'epartements disponibles, la s\'election d'un d\'epartement filtre les mentions, etc. Pour r\'einitialiser, effacer le filtre au niveau sup\'erieur.
\end{notebox}

\subsubsection{Informations affich\'ees dans les r\'esultats}

Pour chaque r\'esultat (responsable trouv\'e), les informations suivantes sont g\'en\'eralement disponibles\,:
\begin{itemize}
  \item Nom et pr\'enom.
  \item R\^ole dans la structure.
  \item Structure d'appartenance (composante, d\'epartement, mention).
  \item Courriel institutionnel.
  \item Courriel fonctionnel (si renseign\'e).
  \item T\'el\'ephone (si renseign\'e).
  \item Bureau (si renseign\'e).
\end{itemize}

\screenshot{img_annuaire_responsables_detail}{Annuaire Responsables -- R\'esultats d\'etaill\'es}{Capturer l'onglet Responsables avec des r\'esultats affich\'es sous forme de cartes ou de tableau. Montrer pour au moins 3 r\'esultats\,: nom/pr\'enom, r\^ole, structure, courriel. Si des informations de t\'el\'ephone ou bureau sont renseign\'ees, les montrer \'egalement.}

\subsection{Onglet Formations}

Cet onglet permet de retrouver les formations et leurs responsables.

\subsubsection{Filtres disponibles}

\begin{itemize}
  \item \textbf{Recherche libre} : nom de la formation.
  \item \textbf{Type de dipl\^ome} : Licence, Master, Doctorat, DUT, BTS, etc.
  \item \textbf{Composante} et \textbf{D\'epartement}.
\end{itemize}

\subsubsection{Informations affich\'ees}

Pour chaque formation\,:
\begin{itemize}
  \item Intitul\'e de la formation.
  \item Type de dipl\^ome.
  \item Composante et d\'epartement.
  \item Directeur de mention / Responsable de formation.
  \item Directeur d'\'etudes (si affect\'e).
\end{itemize}

\screenshot{img_annuaire_formations_detail}{Annuaire Formations -- R\'esultats}{Capturer l'onglet Formations avec des r\'esultats. Montrer les filtres (type de dipl\^ome, composante) et les cartes de r\'esultats avec l'intitul\'e de la formation et les responsables associ\'es.}

\subsection{Onglet Structures}

Cet onglet liste les entit\'es organisationnelles de l'\'etablissement.

\subsubsection{Filtres disponibles}

\begin{itemize}
  \item \textbf{Recherche libre} : nom de la structure.
  \item \textbf{Type de structure} : Composante, D\'epartement, Mention, Parcours, Niveau, Service.
  \item \textbf{Parent} (hi\'erarchie).
\end{itemize}

\subsubsection{Naviguer dans la hi\'erarchie}

\begin{steps}
  \item Dans les r\'esultats, cliquer sur une structure pour ouvrir sa fiche d\'etaill\'ee.
  \item La fiche montre les sous-structures directes (liens cliquables).
  \item Cliquer sur une sous-structure pour naviguer vers elle.
\end{steps}

\screenshot{img_annuaire_structures_detail}{Annuaire Structures -- Navigation hi\'erarchique}{Capturer l'onglet Structures avec des r\'esultats. Montrer plusieurs types de structures (composante, d\'epartement, mention) et une fiche ouverte en p\'eriph\'erie si possible (ou en mode pleine page).}

\subsection{Onglet Secr\'etariats}

Cet onglet affiche les contacts du secr\'etariat p\'edagogique.

\subsubsection{Informations affich\'ees}

Pour chaque agent de secr\'etariat\,:
\begin{itemize}
  \item Nom et pr\'enom.
  \item Structure de rattachement.
  \item Courriel fonctionnel du service.
  \item T\'el\'ephone du secr\'etariat.
  \item Bureau / localisation.
\end{itemize}

\screenshot{img_annuaire_secretariats_detail}{Annuaire Secr\'etariats -- R\'esultats}{Capturer l'onglet Secr\'etariats avec des r\'esultats. Montrer les contacts avec leur structure de rattachement, courriel fonctionnel, t\'el\'ephone et bureau.}

\section{Fiche Structure -- Guide d\'etaill\'e}
\label{sec:fiche_structure}

La fiche structure est accessible en cliquant sur une structure dans les r\'esultats de l'annuaire ou via la rubrique \textbf{Structures}.

\subsection{Sections de la fiche}

\begin{description}[style=nextline, font=\bfseries]
  \item[Identification] Nom, type de structure (Composante, D\'epartement, Mention...), code, ann\'ee acad\'emique.
  \item[Hi\'erarchie] Structure parente (lien cliquable) et niveau dans la hi\'erarchie.
  \item[Coordonn\'ees] Adresse physique, t\'el\'ephone principal, courriel institutionnel.
  \item[Responsables] Liste des personnes affect\'ees avec leur r\^ole, courriel et coordonn\'ees.
  \item[Secr\'etariat] Liste des agents de secr\'etariat avec leurs contacts fonctionnels.
  \item[Sous-structures] Liens vers les entit\'es filles directes (d\'epartements, mentions...).
  \item[Statistiques] Nombre de responsables, d'affectations, de sous-structures.
\end{description}

\screenshot{img_fiche_structure_complete}{Fiche structure compl\`ete -- Toutes sections}{Capturer une fiche structure enti\`ere (plusieurs d\'efilements si n\'ecessaire). Montrer successivement toutes les sections\,: identification, hi\'erarchie, coordonn\'ees, responsables, secr\'etariat, sous-structures. Utiliser une composante pour avoir le maximum de contenu.}

\section{Organigrammes -- Guide de consultation}
\label{sec:organigrammes_consultation}

\subsection{Vue Structures}

La vue Structures repr\'esente l'organisation hi\'erarchique des entit\'es.

\begin{itemize}
  \item Chaque n\oe ud affiche\,: nom de la structure, type, responsable(s) principal(aux).
  \item Les li\`egnes repr\'esentent les relations parent-enfant dans la hi\'erarchie.
  \item Il est possible de d\'eplier/replier les branches.
  \item Un clic sur un n\oe ud peut afficher les d\'etails de la structure.
\end{itemize}

\screenshot{img_organigramme_vue_structures}{Organigramme -- Vue Structures d\'etaill\'ee}{Capturer l'organigramme en vue Structures avec plusieurs niveaux d\'epli\'es (composante $\to$ d\'epartements $\to$ mentions si possible). Montrer quelques n\oe uds avec leurs informations visibles (nom structure, type, responsable).}

\subsection{Vue Personnes}

La vue Personnes repr\'esente les liens hi\'erarchiques N+1 entre personnes.

\begin{itemize}
  \item Chaque n\oe ud affiche\,: nom, pr\'enom, r\^ole, structure d'affiliation, courriel.
  \item Les liens repr\'esentent la relation de supervision (N+1).
  \item Les filtres permettent de pr\'eciser\,: r\^ole, composante, d\'epartement, mention, parcours, niveau.
\end{itemize}

\screenshot{img_organigramme_vue_personnes}{Organigramme -- Vue Personnes d\'etaill\'ee}{Capturer l'organigramme en vue Personnes avec les filtres visibles et plusieurs n\oe uds personnes affich\'es. Un n\oe ud doit \^etre survol\'e ou ouvert pour montrer les informations disponibles (r\^ole, structure, courriel).}

\subsection{Filtres de l'organigramme}

Les filtres disponibles dans la vue Personnes\,:

\medskip
\begin{tabularx}{\textwidth}{L{4cm}X}
  \toprule
  \rowcolor{cgplightblue}
  \textbf{Filtre} & \textbf{Effet} \\
  \midrule
  Recherche libre  & Filtre les n\oe uds par nom de personne ou de structure. \\
  \rowcolor{cgplightgray}
  R\^ole           & N'affiche que les personnes ayant ce r\^ole. \\
  Composante       & Limite l'arbre \`a une composante. \\
  \rowcolor{cgplightgray}
  D\'epartement    & Limite l'arbre \`a un d\'epartement. \\
  Mention          & Limite l'arbre \`a une mention. \\
  \rowcolor{cgplightgray}
  Parcours         & Limite \`a un parcours. \\
  Niveau           & Limite \`a un niveau d'ann\'ee. \\
  \bottomrule
\end{tabularx}

\subsection{Exporter un organigramme}

Les utilisateurs autoris\'es (voir chapitre~\ref{chap:roles}) peuvent exporter les organigrammes affich\'es\,:

\begin{steps}
  \item Afficher l'organigramme souhait\'e.
  \item Cliquer sur le bouton \textbf{Exporter}.
  \item S\'electionner le format\,: \textbf{PDF} (pour impression), \textbf{SVG} (vectoriel), \textbf{CSV} (tabulaire), \textbf{JSON} (donn\'ees brutes).
\end{steps}

\screenshot{img_organigramme_export_formats}{Export d'organigramme -- S\'election du format}{Capturer le menu ou la modale d'export affichant les formats disponibles (PDF, SVG, CSV, JSON).}

\section{Signalements -- Guide d\'etaill\'e}
\label{sec:signalements_detail}

\subsection{Types de signalements}

\medskip
\begin{tabularx}{\textwidth}{L{3cm}X}
  \toprule
  \rowcolor{cgplightblue}
  \textbf{Type} & \textbf{Utilisation} \\
  \midrule
  Nom            & Erreur sur le nom de famille d'une personne. \\
  \rowcolor{cgplightgray}
  Pr\'enom       & Erreur sur le pr\'enom. \\
  T\'el\'ephone  & Num\'ero de t\'el\'ephone incorrect ou manquant. \\
  \rowcolor{cgplightgray}
  Bureau         & Localisation (bureau, b\^atiment) incorrecte. \\
  Courriel       & Adresse \'electronique incorrecte ou manquante. \\
  \rowcolor{cgplightgray}
  R\^ole         & Affectation manquante, incorrecte ou obsol\`ete. \\
  Autre          & Tout autre type de probl\`eme sur les donn\'ees. \\
  \bottomrule
\end{tabularx}

\subsection{Cycle de vie d'un signalement}

\begin{center}
\begin{tikzpicture}[
  node distance=1.5cm,
  state/.style={rectangle, rounded corners=6pt, draw=cgpblue, fill=cgplightblue,
                text width=3cm, align=center, font=\small\bfseries, minimum height=0.9cm}
]
  \node[state] (soumis) {Soumis\\(En attente)};
  \node[state, right=2.5cm of soumis] (encours) {En cours\\de traitement};
  \node[state, right=2.5cm of encours] (clos) {Cl\^otur\'e};
  \node[state, below=1.2cm of encours] (escalade) {Escalad\'e\\(SC)};

  \draw[->, thick, cgpblue] (soumis) -- (encours)
    node[midway, above, font=\tiny] {Pris en charge};
  \draw[->, thick, cgpblue] (encours) -- (clos)
    node[midway, above, font=\tiny] {R\'esolu};
  \draw[->, thick, cgpblue] (encours) -- (escalade)
    node[midway, right, font=\tiny] {Escalade};
  \draw[->, thick, cgpblue] (escalade) -- (clos)
    node[midway, below right, font=\tiny] {R\'esolu par SC};
\end{tikzpicture}
\end{center}

\subsection{Soumettre un signalement}

\begin{steps}
  \item Naviguer vers \textbf{Signalements}.
  \item Cliquer sur \textbf{Cr\'eer un signalement}.
  \item S\'electionner le \textit{type d'erreur}.
  \item Rechercher et s\'electionner la \textit{personne} ou la \textit{structure} concern\'ee.
  \item Saisir une \textit{description} pr\'ecise\,: d\'ecrire l'erreur constat\'ee et si possible indiquer la valeur correcte.
  \item Cliquer sur \textbf{Envoyer}.
\end{steps}

\screenshot{img_signalement_formulaire}{Formulaire de signalement}{Capturer le formulaire de cr\'eation d'un signalement avec tous les champs remplis\,: menu d\'eroulant du type d'erreur (montrer les options), champ de recherche de la personne/structure concern\'ee avec un exemple de r\'esultat s\'electionn\'e, et le champ description avec un texte exemple.}

\subsection{Suivre ses signalements}

\begin{steps}
  \item Naviguer vers \textbf{Signalements}.
  \item La liste affiche les signalements soumis avec leur statut.
  \item Cliquer sur un signalement pour voir ses d\'etails et les \'eventuels commentaires de traitement.
\end{steps}

\screenshot{img_signalements_liste}{Liste des signalements avec statuts}{Capturer la liste des signalements avec les badges de statut color\'es (En attente = orange, En cours = bleu, Cl\^otur\'e = vert). Montrer plusieurs signalements avec des statuts diff\'erents.}

\section{Mon Profil -- Guide d\'etaill\'e}
\label{sec:profil_detail}

La page \textbf{Mon Profil} est accessible depuis le menu ou l'ic\^one utilisateur.

\subsection{Informations consultables}

\begin{itemize}
  \item \textbf{Identit\'e} : nom, pr\'enom (informations de r\'ef\'erence, non modifiables par l'utilisateur).
  \item \textbf{Identifiant de connexion} (non modifiable).
  \item \textbf{Ann\'ee courante} et \textbf{r\^oles actifs} : liste de toutes les affectations pour l'ann\'ee en cours, avec la structure et le r\^ole correspondants.
  \item \textbf{Coordonn\'ees personnelles} : courriel secondaire, t\'el\'ephone, bureau.
\end{itemize}

\subsection{Informations modifiables}

\begin{warnbox}[R\`egle de modification du profil]
  L'utilisateur peut uniquement modifier ses \textbf{coordonn\'ees personnelles} (courriel secondaire, t\'el\'ephone, bureau). \\
  Les informations d'identit\'e (nom, pr\'enom) et les affectations ne peuvent \^etre modifi\'ees que par un administrateur ou les Services Centraux. En cas d'erreur, cr\'eer un signalement.
\end{warnbox}

\begin{steps}
  \item Naviguer vers \textbf{Mon Profil}.
  \item Cliquer sur \textbf{Modifier} dans la section \og Mes coordonn\'ees \fg{}.
  \item Mettre \`a jour les champs souhait\'es\,:
    \begin{itemize}
      \item \textit{Courriel secondaire} : adresse \'electronique personnelle ou alternative.
      \item \textit{T\'el\'ephone} : num\'ero de t\'el\'ephone direct.
      \item \textit{Bureau} : num\'ero de bureau, b\^atiment, campus.
    \end{itemize}
  \item Cliquer sur \textbf{Enregistrer}.
\end{steps}

\screenshot{img_profil_modification}{Profil -- Mode modification}{Capturer la page Mon Profil en mode modification (apr\`es clic sur Modifier). Montrer les champs \'editables (courriel secondaire, t\'el\'ephone, bureau) avec leurs valeurs actuelles, et les boutons Enregistrer et Annuler.}

\section{Notifications -- Guide d\'etaill\'e}
\label{sec:notifications_detail}

Le syst\`eme de notifications informe les utilisateurs des \'ev\'enements les concernant.

\subsection{Types de notifications}

\medskip
\begin{tabularx}{\textwidth}{L{5cm}X}
  \toprule
  \rowcolor{cgplightblue}
  \textbf{\'Ev\'enement} & \textbf{Destinataire} \\
  \midrule
  Demande de r\^ole approuv\'ee   & Demandeur \\
  \rowcolor{cgplightgray}
  Demande de r\^ole refus\'ee     & Demandeur \\
  Signalement cl\^otur\'e        & Auteur du signalement \\
  \rowcolor{cgplightgray}
  Signalement escalad\'e         & Services Centraux \\
  D\'el\'egation re\c{c}ue       & D\'el\'egu\'e \\
  \rowcolor{cgplightgray}
  D\'el\'egation r\'evoqu\'ee    & D\'el\'egu\'e \\
  \bottomrule
\end{tabularx}

\subsection{Consulter les notifications}

\begin{steps}
  \item Cliquer sur la cloche \faBell{} dans l'en-t\^ete (le badge indique le nombre de non-lues).
  \item Le panneau de notifications s'ouvre.
  \item Cliquer sur une notification pour la marquer comme lue et acc\'eder \`a la ressource concern\'ee.
  \item Optionnel\,: cliquer sur \textbf{Tout marquer comme lu} pour effacer toutes les notifications.
\end{steps}

\screenshot{img_notifications_detail}{Panneau notifications ouvert}{Capturer le panneau de notifications ouvert avec plusieurs notifications de types diff\'erents (demande approuv\'ee, signalement cl\^otur\'e, d\'el\'egation). Montrer les notifications lues et non lues avec une distinction visuelle claire.}

\section{Raccourcis clavier et navigabilit\'e}

\begin{tipbox}[Navigation rapide]
  \begin{itemize}
    \item \textbf{Tabulation} : naviguer entre les champs de formulaire.
    \item \textbf{Entr\'ee} : valider un formulaire ou une recherche.
    \item \textbf{\'Echap} : fermer une modale ou annuler une action.
    \item \textbf{Ctrl+F / Cmd+F} : recherche dans la page (navigateur).
  \end{itemize}
\end{tipbox}

\section{Accessibilit\'e et bonnes pratiques}

\begin{itemize}
  \item Utiliser de pr\'ef\'erence \textbf{Google Chrome} ou \textbf{Firefox} en version r\'ecente.
  \item \'Eviter d'ouvrir l'application dans plusieurs onglets simultan\'ement pour \'eviter les conflits de session.
  \item En cas de comportement inattendu, rafra\^ichir la page (\textsf{F5}) avant de signaler un bug.
  \item Ne pas partager vos identifiants de connexion.
  \item Si vous quittez votre poste de travail, \textbf{d\'econnectez-vous} de l'application.
\end{itemize}


\end{document}
